\FloatBarrier
\newpage
\section{Comparative pore analysis using physical and imaging techniques}

The determination of porosity is fundamental to understanding and predicting the performance and durability of concrete.
Whilst macroscopic air voids influence properties like freeze-thaw resistance, the gel porosity dictates the material's intrinsic characteristics.
For example, the transport mechanisms (permeability and diffusion) of water and ions are determined by the pore network which is formed, among other things, by the fine structures inherent of the \ac{CSH} phase.
Accurately quantifying the gel and the capillary porosity is essential to assess the resistance to chemical (e.g. chloride ingress or sulphate attack) or frost attack, its mechanical performance and durability.
This becomes especially crucial for modern composite cements, with a much more complex microstructure.

On a macroscopic scale, comparing different pore measurement methods is feasible.
However, as the pore size decreases, this comparison becomes increasingly challenging due to physical limits and the fractal nature of porosity.
To evaluate and compare the results of these methods, a standard material was utilised to benchmark and compare selected methods, which were described in \zcref{sec:methods_pore_analysis}.

\subsection{Porous glass standard}

Two porous glass standards (\textsc{Vycor} glass with \num{5} and \qty{20}{\nano\meter} pore diameter, short: \acs{VG5} and \acs{VG20}, see \zcref{sec:boraglas}) were evaluated using \ac{SEM} and \ac{LTDSC} spectroscopy.
The glass standards were initially defined using \ac{MIP} by the manufacturer (\zcref{fig:vycor-mip}). % or \ce{N2}-ad/absorption.
% TODO: die Endporosität wird normalerweise eher durch Engstellen (necks) definiert als durch die Bulk-porosität

\subsubsection{\acs*{SEM}}

To obtain \ac{SEM} images, the glass was fractured and even surfaces were imaged without any coating \cite{kleiner.2026i}.
As the \ac{SEM} images show (\zcref{fig:ltdsc-vycor-glass-a, fig:ltdsc-vycor-glass-b}), the pores are interconnected, which was also already demonstrated in literature \cite{levitz.1998, levitz.2007}.
This pore structure allows the use of this glass as a filter material and is caused by the manufacturing process used (sintering of glass powder).

% OLD: C:\Users\Florian Kleiner\Nextcloud\Uni\WdB\Paper\2024 LT-DSC-XRD-BSE Alit-Belit Paper\5. REM+EDX\Nanoporöses Glas\2022_02_08 Poröses Glas
% NEW: X:\Mitarbeiter\Florian Kleiner\Einzelbilder\2025 Vycor Auswertung\
% 3 * 5.18x3.45 µm² = 3 * 17.871 µm² = 53.61 µm²
\begin{figure}[htb!]
    \centering

	\begin{subfigure}{0.49\textwidth}
		\includegraphics[width=\textwidth]{ltdsc/glass/segmentation/5nm.pdf}
		\caption{Raw, \acs{VG5}}
		\label{fig:ltdsc-vycor-glass-a}
	\end{subfigure}
	\hfill
	\begin{subfigure}{0.49\textwidth}
		\includegraphics[width=\textwidth]{ltdsc/glass/segmentation/20nm.pdf}
		\caption{Raw, \acs{VG20}}
		\label{fig:ltdsc-vycor-glass-b}
	\end{subfigure}

	\begin{subfigure}{0.49\textwidth}
		\includegraphics[width=\textwidth]{ltdsc/glass/segmentation/5nm_pores}
		\caption{Segmented, \acs{VG5}}
		\label{fig:ltdsc-vycor-glass-c}
	\end{subfigure}
	\hfill
	\begin{subfigure}{0.49\textwidth}
		\includegraphics[width=\textwidth]{ltdsc/glass/segmentation/20nm_pores.pdf}
		\caption{Segmented, \acs{VG20}}
		\label{fig:ltdsc-vycor-glass-d}
	\end{subfigure}

	\begin{subfigure}{0.49\textwidth}
		\includegraphics[width=\textwidth]{ltdsc/glass/segmentation/5nm_pores_ws}
		\caption{Segmented, watershed filter, \acs{VG5}}
		\label{fig:ltdsc-vycor-glass-e}
	\end{subfigure}
	\hfill
	\begin{subfigure}{0.49\textwidth}
		\includegraphics[width=\textwidth]{ltdsc/glass/segmentation/20nm_pores_ws.pdf}
		\caption{Segmented, watershed filter, \acs{VG20}}
		\label{fig:ltdsc-vycor-glass-f}
	\end{subfigure}
    \caption{
		Enlarged sections of \ac{SEM}-\ac{SE} images of two \textsc{Vycor}-glass samples and their respective segmentation results \cite{kleiner.2026i}.
		The pore size was determined using two segmented \ac{SEM} images for each glass type.
		In (e) and (f), each pore was assigned a random colour for visualisation purposes.
	} % Probenbilder-rohdaten vmtl vertauscht?
    \label{fig:ltdsc-vycor-glass}
\end{figure}


Since the dwell-time has to be chosen as short as possible to avoid charging, the images became very noisy and were denoised using a median blur filter using a \numproduct{3x3} kernel.

The pores were segmented using adaptive thresholding to determine the pore size distribution (see \zcref{fig:ltdsc-vycor-glass-c,fig:ltdsc-vycor-glass-d}).
This method utilises a `blocksize' parameter (similar to a sliding window; in \unit{\px}), which defines the local neighbourhood size for calculating the threshold.
For \acs{VG20}, a blocksize corresponding to \qtyrange{50}{100}{\nano\meter} yielded reliable results.
Consequently, a blocksize parameter corresponding to approximately \qty{80}{\nano\meter} was selected.
The resulting binary images were morphologically cleaned using dilation and erosion operations.

%As the glass is manufactured as a filter material by sintering glass powder, a tunnel-like pore network is formed.
Since the pores in the glass are interconnected, the segmented pores in the \ac{SEM}-\ac{SE} images often appear larger.
This can lead to a significant overestimation of pore sizes.
To address this issue, a watershed-based separation algorithm was applied for comparison.
The minimum spacing between the maxima was set to \qty{3.5}{\nano\meter}.
This parameter was selected manually as it resulted in a subjectively adequate separation of pores into fairly compact shapes (compare \zcref{fig:ltdsc-vycor-glass-e, fig:ltdsc-vycor-glass-f}).
However, it should be noted that variations in this parameter significantly influence the final result.

The resulting \ac{PSD}s, derived from three images per specimen, are presented in \zcref{fig:compare_psd_vycor}.
This figure illustrates both the raw segmentation (blue) and the influence of the watershed algorithm (`ws', red).
To determine the dominant pore radius, the histograms were fitted using \zcref{eq:weig_norm_pore_frequency}.
The fit parameters are listed in \zcref{tab:rP_fit_parameters}.


\begin{figure}[htbp]
    \centering
    % 5nm Graph
    \begin{subfigure}{0.49\textwidth}
        \flushright
        \begin{tikzpicture}
            \begin{axis}[
                width=1.1\linewidth,
                height=6cm,
                grid=major,
                xlabel={pore radius in \unit{\nano\meter}},
                ylabel={pore frequency $f_\text{v,n}$},
				ytick=\empty,
                xmin=0, xmax=54,
                ymin=0, ymax=1.2,
                grid style={line width=.1pt, draw=gray!10},
                legend style={at={(0.98,0.98)},anchor=north east, font=\scriptsize},
                legend columns=1,
            ]

                \addplot[ybar, bar width=1.6, fill opacity=0.5, draw=none, area legend, fill=blue] table[x=bin_center,y expr=\thisrow{count_per_um2}/510161.7,col sep=comma] {figures/ltdsc/glass/segmentation/5nm_PSD_volume.csv};
                \addlegendentry{\acs*{PSD}};

                \addplot[ybar, bar width=1.6, fill opacity=0.5, draw=none, area legend, fill=red] table[x=bin_center,y expr=\thisrow{count_per_um2}/663987.5,col sep=comma] {figures/ltdsc/glass/segmentation/5nm_PSD_volume_ws.csv};
                \addlegendentry{\acs*{PSD} ws};

				\addplot [no markers, blue, domain=0.01:55, samples=100] {(6861582.0403 / x * exp(-0.5 * ((ln(x) - 2.6318) / 0.3063)^2))/510161.7};
				\addlegendentry{\acs*{PSD} fit}
				% max value at x=12.6538

				\addplot [no markers, red, domain=0.01:55, samples=100] {(5695958.6516 / x * exp(-0.5 * ((ln(x) - 2.1651) / 0.2530)^2)) / 663987.5};
				\addlegendentry{\acs*{PSD} ws fit}
				% max value at x=8.1751

				%--- Fit Result for PSD_volume_weighted.png ---
				% Peak: 517409.3924 at x=12.6538; max at
				% \addplot [no markers, red, domain=0.01:55, samples=100] {6861582.0403 / x * exp(-0.5 * ((ln(x) - 2.6318) / 0.3063)^2)};


            \end{axis}
        \end{tikzpicture}
        \caption{\acs{VG5}.}
        \label{fig:compare_psd_5nm}
    \end{subfigure}
    \hfill
    % 20nm Graph
    \begin{subfigure}{0.49\textwidth}
        \flushleft
        \begin{tikzpicture}
            \begin{axis}[
                width=1.1\linewidth,
                height=6cm,
                grid=major,
                xlabel={pore radius in \unit{\nano\meter}},
                ylabel={pore frequency $f_\text{v,n}$},
                yticklabel pos=right,
				ytick=\empty,
                xmin=0, xmax=54,
                ymin=0, ymax=1.2,
                grid style={line width=.1pt, draw=gray!10},
                legend style={at={(0.98,0.98)},anchor=north east, font=\scriptsize},
            ]

                \addplot[ybar, bar width=1.6, fill opacity=0.5, draw=none, area legend, fill=blue] table[x=bin_center,y expr=\thisrow{count_per_um2}/597423.3,col sep=comma] {figures/ltdsc/glass/segmentation/20nm_PSD_volume.csv};
                \addlegendentry{\acs*{PSD}};

                \addplot[ybar, bar width=1.6, fill opacity=0.5, draw=none, area legend, fill=red] table[x=bin_center,y expr=\thisrow{count_per_um2}/791589.7,col sep=comma] {figures/ltdsc/glass/segmentation/20nm_PSD_volume_ws.csv};
                \addlegendentry{\acs*{PSD} ws};

				\addplot [no markers, blue, domain=0.01:55, samples=100] {(12245826.5777 / x * exp(-0.5 * ((ln(x) - 3.1318) / 0.3456)^2)) / 597423.3};
				\addlegendentry{\acs*{PSD} fit}
				% max value at x=20.3357

				\addplot [no markers, red, domain=0.01:55, samples=100] {(9220842.0696 / x * exp(-0.5 * ((ln(x) - 2.4632) / 0.2785)^2))/791589.7};
				\addlegendentry{\acs*{PSD} ws fit}
				% max value at x=10.8658

            \end{axis}
        \end{tikzpicture}
        \caption{\acs{VG20}.}
        \label{fig:compare_psd_20nm}
    \end{subfigure}
	\caption{
		\acs*{PSD} of \acs{VG5} and \acs{VG20}, expressed as the volume-weighted and normalised pore frequency $f_\text{v,n}$.
		The distributions were obtained from segmented \ac{SE} images (\num{3} images, \qty{53.61}{\micro\meter\squared} in total for each glass).
		Graphs marked with `ws' were filtered using the watershed algorithm \cite{beucher.1992}.
		The fit parameters are listed in \zcref{tab:rP_fit_parameters}.
	}
    \label{fig:compare_psd_vycor}
\end{figure}


For \acs{VG5}, a main radius of approximately \qty{12.7}{\nano\meter} was calculated.
Applying the watershed filter reduced this radius to \qty{8.2}{\nano\meter}.
For \acs{VG20}, radii of \num{20.3} and \qty{10.9}{\nano\meter} were determined, respectively.
In both cases, the calculated main radii overestimate the manufacturer's specifications.
For \acs{VG5}, the calculated values exceed the manufacturer's median pore radius of \qty{2.4}{\nano\meter} by a factor of \numrange{3.5}{5.4}.
For \acs{VG20}, the overestimation is lower, exceeding the reference radius of \qty{9.2}{\nano\meter} (measured via \ac{MIP}, see \zcref{fig:vycor-mip}) by a factor of \numrange{1.2}{2.2}.

% TODO: Median pore size eintragen ins Diagram
% TODO: normales thresholding auf eines der identischen Bilder anwenden zum Vergleich?


\FloatBarrier
\subsubsection{\acs*{LTDSC}}

The glass specimens were measured using \ac{LTDSC} \cite{kleiner.2026i} utilising a simplified temperature regimen compared to the regime presented in \zcref{fig:ltdc_regime}.
It also experienced ice crystal-inducing freeze-thaw cycles before the actual measurement.
For the measurement, the water saturated specimen was then cooled down to \qty{-60}{\celsius} and heated up to \qty{20}{\celsius} at a cooling and heating rate of $\pm$\,\qty{2}{\kelvin\per\minute}.

\begin{figure}[!htb]
    \centering

 	\begin{subfigure}{.48\textwidth}
      \begin{tikzpicture}
        \begin{semilogxaxis}[
				blue,
				axis line style={black},
				xmin=1.5, xmax=60,
				ymin=-.20, ymax=0,
				y dir=reverse,
				width=\linewidth,
				height=6cm,
				grid=none,
				xlabel={temperature in \unit{\celsius}},
				xticklabels={\num{-51.2}, \num{-25.3}, \num{-12.4}, \num{-4.6}, \num{-2.0}},
				xtick={2.5,5,10,25,50},
				ytick=\empty,
				scaled y ticks=false,
				scaled x ticks=false,
				axis x line*=top
            ]

        \end{semilogxaxis}
        \begin{semilogxaxis}[
				legend style={
					at={(0.015,.8)},
					anchor=west,
					cells={align=left},
					font=\footnotesize
				},
				xmin=1.5, xmax=60,
				ymin=-.20, ymax=0,
				y dir=reverse,
				width=\linewidth,
				height=6cm,
				grid=major,
				xlabel={pore radius $r_\text{P}$ in \unit{\nano\meter}},
				ylabel={heat flow $\dot{Q}$ in \unit{\joule\per\second\per\gram}},
				grid style={line width=.1pt, draw=gray!10},
				yticklabel style={
						/pgf/number format/fixed,
						/pgf/number format/precision=3
				},
				xticklabels={2.5,5,10,25,50},
				xtick={2.5,5,10,25,50},
				ytick={0,-0.1,-0.2},
				scaled y ticks=false,
				scaled x ticks=false,
            ]
			  \addplot[dotted, thick, samples=2, darkgrey,forget plot] coordinates {(3.19,-.2)(3.19,.2)};
			  %\node[darkgrey, rotate=90, anchor=west] at (axis cs: 2.6,-.06) {\scriptsize <\,\qty{3.19}{\nano\meter}}
			  \node[darkgrey, rotate=90, anchor=west] at (axis cs: 2.6,-.06) {\scriptsize <\,\qty{3.19}{\nano\meter}};

			\addplot [ no marks, darkgreen] table [each nth point=4, filter discard warning=false, unbounded coords=discard,x index=2, y index=1, col sep=comma] {figures/ltdsc/glass/ExpDat_2.1-5nm_freeze_segment_18.csv};
			\addlegendentry{specimen 1};
			\addplot [ no marks, darkgreen, dashed] table [each nth point=4, filter discard warning=false, unbounded coords=discard,x index=2, y index=1, col sep=comma] {figures/ltdsc/glass/ExpDat_2.2-5nm_freeze_segment_18.csv};
			\addlegendentry{specimen 2};

        \end{semilogxaxis}
	  \end{tikzpicture}
	  \caption{\acs{VG5}.}
    	\label{fig:ltdsc_result_glass_freeze_a}

	\end{subfigure}
	\hfill
	\begin{subfigure}{.48\textwidth}

		\begin{tikzpicture}
			\begin{semilogxaxis}[
				blue,
				axis line style={black},
					xmin=1.5, xmax=60,
					ymin=-.20, ymax=0,
					y dir=reverse,
					width=\linewidth,
					height=6cm,
					grid=none,
					xlabel={temperature in \unit{\celsius}},
					xticklabels={\num{-51.2}, \num{-25.3}, \num{-12.4}, \num{-4.6}, \num{-2.0}},
					xtick={2.5,5,10,25,50},
					ytick=\empty,
					scaled y ticks=false,
					scaled x ticks=false,
					axis x line*=top,
				]

		  \end{semilogxaxis}
		  \begin{semilogxaxis}[
				legend style={
					at={(0.015,.8)},
					anchor=west,
					cells={align=left},
					font=\footnotesize
				},
				xmin=1.5, xmax=60,
				ymin=-.20, ymax=0,
				y dir=reverse,
				width=\linewidth,
				height=6cm,
				grid=major,
				xlabel={pore radius $r_\text{P}$ in \unit{\nano\meter}},
				ylabel={heat flow $\dot{Q}$ in \unit{\joule\per\second\per\gram}},
				grid style={line width=.1pt, draw=gray!10},
				yticklabel style={
						/pgf/number format/fixed,
						/pgf/number format/precision=3
				},
				xticklabels={2.5,5,10,25,50},
				xtick={2.5,5,10,25,50},
				ytick={0,-0.1,-0.2},
				scaled y ticks=false,
				scaled x ticks=false,
				yticklabel pos=right,
			]
			  \addplot[dotted, thick, samples=2, darkgrey,forget plot] coordinates {(3.19,-.2)(3.19,0)};
			  %\node[darkgrey, rotate=90, anchor=west] at (axis cs: 2.6,-.06) {\scriptsize <\,\qty{3.19}{\nano\meter}}
			  \node[darkgrey, rotate=90, anchor=west] at (axis cs: 2.6,-.06) {\scriptsize <\,\qty{3.19}{\nano\meter}};

			  \addplot [ no marks, orange ] table [ each nth point=4, filter discard warning=false, unbounded coords=discard, x index=2, y index=1, col sep=comma ] {figures/ltdsc/glass/ExpDat_1.1-20nm_freeze_segment_18.csv};
			  \addlegendentry{specimen 1};
			  \addplot [ no marks, orange, dashed ] table [ each nth point=4, filter discard warning=false, unbounded coords=discard, x index=2, y index=1, col sep=comma ] {figures/ltdsc/glass/ExpDat_1.2-20nm_freeze_segment_18.csv};
			  \addlegendentry{specimen 2};

		  \end{semilogxaxis}
		\end{tikzpicture}
		\caption{\acs{VG20}.}
    	\label{fig:ltdsc_result_glass_freeze_b}

	\end{subfigure}
    \caption{
		\ac{LTDSC}-measurements of two nano-porous glass types \cite{kleiner.2026i}.
		Each type was measured using two specimens.
		Note that pore sizes below \qty{3.19}{\nano\meter} (dotted line) are undefined as noted in \zcref{eq:brun_freeze}.
	}
    \label{fig:ltdsc_result_glass_freeze}
\end{figure}

This resulted in the following freeze curves as plotted in \zcref{fig:ltdsc_result_glass_freeze}.
\acs{VG5} specified with a pore radius of \qty{2.5}{\nano\meter} has two peaks in the heat-flow signal (\zcref{fig:ltdsc_result_glass_freeze_a}).
According to \zcref{eq:brun_freeze} \cite{brun.1977}, a heat flow signal was measured in the pore radius range of \qtyrange{3}{8}{\nano\meter} (pore diameter of \num{6} and \qty{16}{\nano\meter}) with peaks at \num{4} and \qty{6.5}{\nano\meter} (pore diameter of \num{8} and \qty{13}{\nano\meter}).

In contrast, \acs{VG20} (\zcref{fig:ltdsc_result_glass_freeze_a}) has a specified pore radius of \qty{10}{\nano\meter}. The measured heat flow signal can be converted to pore radii of \num{5} and \qty{12}{\nano\meter}, with peaks at pore radii of \num{7} and \qty{10}{\nano\meter}.

These results show that the pore radius is overestimated by a factor of up to \num{2} for \acs{VG5} compared to the \ac{MIP} results (see \zcref{fig:vycor-mip}), while \acs{VG20} shows a good fit for its main peak.
This result is similar to the image based results with the applied watershed filter shown before.
Despite this deviation, the difference in the pore size is still observable.
Compared to the \ac{SEM} results (see \zcref{fig:ltdsc-vycor-glass}), the \ac{LTDSC} results fit slightly better to the \ac{MIP} results.
Similar to the \ac{SEM} data, the \qty{5}{\nano\meter} glass also exhibits a shift towards larger pore sizes compared to the \ac{MIP} results.



\FloatBarrier
\subsection{Porosity of hydrated \acs*{C2S} and \acs*{C3S}}

In contrast to porous glass, the generation of fractured surfaces in hydrated \ac{C2S} and \ac{C3S} is not sufficient for a pore analysis using \ac{SEM}-\ac{BSE} images.
Therefore, a good specimen preparation is important to obtain a suitable pore analysis and \ac{ArBIB} sectioned specimens were used (see \zcref{sec:bib_pol_fixed_results}).

\subsubsection{\acs*{SEM}}
\label{subsec:pores_C3S_SEM}

\textit{This section was partially published in \cite{kleiner.2021a}, but the dataset was reevaluated.}

The images obtained from the \ac{ArBIB} sectioned specimens (see \zcref{fig:csh-arbib-SEM2, fig:csh-arbib-SEM}) reveal the porosity within \ac{CSHi} and the porosity caused by the interpenetrating network of \ac{CSH} fibres.
This also illustrates the advantage of the epoxy resin intrusion (\zcref{fig:csh-arbib-b}) into the pore space between the needle structure for image segmentation.
In \zcref{fig:csh-arbib-a, fig:csh-arbib-c}, the full pore background is visible, making a useful segmentation and evaluation of the pore space unfeasible.
However, the image quality of the nano-scale porosity within \ac{CSHi} is not affected by the epoxy intrusion.
Therefore, a segmentation proposed in \zcref{subsec:algo_psd} can be applied for both preparation variants.

To assess potential damage induced by thermal stress during the \ac{ArBIB} process at room temperature, the results were compared with those obtained using cryo-\ac{ArBIB}.
As no subjective alteration of the \ac{CSH} structures was observed in high-resolution images following either preparation method, a pore analysis of the \ac{CSHi} was conducted.
A total of \num{11} images (total image area: \qty{29.8}{\micro\meter\squared}, of which \qty{23.2}{\micro\meter\squared} was \ac{CSHi}) were analysed for specimens prepared at \qty{20}{\celsius}, and \num{15} images (total image area: \qty{25.2}{\micro\meter\squared}, of which \qty{24.0}{\micro\meter\squared} was \ac{CSHi}) for those prepared at \qty{-140}{\celsius}.
These images, along with the corresponding segmentation results, are provided in \zcref{subsec:appendix_PSD_arbib} and \cite{kleiner.2021g,kleiner.2021h}.
\zcref{fig:compare_images} presents two example images alongside their respective pore segmentations.
Additional image examples, indicating the excluded areas, are shown in \zcref{fig:compare_rt_1,fig:compare_cryo_1} in the Appendix.

\begin{figure}[htb!]
    \centering
    % RT Images
	% C3S 28d BIB 6KV 6h_019.tif
    \begin{subfigure}{0.49\textwidth}
        \includegraphics[width=\textwidth]{methods/arbib-c3s-pores/C3S 28d BIB 6KV 6h_019-1.pdf}
        \caption{\ac{SEM}-\ac{SE} image, \ac{ArBIB} at \qty{20}{\celsius}.}
        \label{fig:compare_rt_a}
    \end{subfigure}
    \hfill
    % Cryo Images
	% C3S 28d cryo BIB 6h_016.tif
    \begin{subfigure}{0.49\textwidth}
        \includegraphics[width=\textwidth]{methods/arbib-c3s-pores/C3S 28d cryo BIB 6h_015-1.pdf}
        \caption{\ac{SEM}-\ac{SE} image, \ac{ArBIB} at \qty{-140}{\celsius}.}
        \label{fig:compare_cryo_a}
    \end{subfigure}

    \begin{subfigure}{0.49\textwidth}
        \includegraphics[width=\textwidth]{methods/arbib-c3s-pores/C3S 28d BIB 6KV 6h_019-1_pores.pdf}
        \caption{Segmented pores, \ac{ArBIB} at \qty{20}{\celsius}.}
        \label{fig:compare_rt_b}
    \end{subfigure}
    \hfill
    \begin{subfigure}{0.49\textwidth}
        \includegraphics[width=\textwidth]{methods/arbib-c3s-pores/C3S 28d cryo BIB 6h_015-1_pores.pdf}
        \caption{Segmented pores, \ac{ArBIB} at \qty{-140}{\celsius}.}
        \label{fig:compare_cryo_b}
    \end{subfigure}
    \caption{
		Comparison of contrast enhanced \ac{SEM}-\ac{SE} images (\qtyproduct{2022 x 1650}{\px} at \qty{0.9}{\nano\meter\per\px}) and the segmentation of pores within \ac{CSHi}, sectioned using \ac{ArBIB} prepared at \num{20} and \qty{-140}{\celsius}.
	}
    \label{fig:compare_images}
\end{figure}

The \ac{PSD} and \ac{CDF} were calculated for each image.
As the pixel sizes and image areas differed, these results subsequently summarised in two datasets.
For each dataset, the area-weighted \ac{PSD} and the corresponding \ac{CDF} were calculated using \zcref{eq:weig_norm_pore_frequency} and plotted in \zcref{fig:compare_graphs}.

\zcref{fig:compare_images} shows nano-sized pores in a sample sectioned using \ac{ArBIB} at \qty{-140}{\celsius} and \qty{20}{\celsius}.
Pores with a diameter of \qty{5}{\nano\meter} can be clearly distinguished and analysed in the \ac{SE} images.
Manual verification of the segmentation results is essential to exclude erroneously segmented areas, as the false detection of microstructural features remains a possibility.

\begin{figure}[htb!]
    \centering
    % RT Graph
    \begin{subfigure}{0.49\textwidth}
        \flushright
        \begin{tikzpicture}
            \begin{axis}[
                width=1.1\linewidth,
                height=6cm,
                grid=major,
                xlabel={pore radius in \unit{\nano\meter}},
                ylabel={weighted frequency $f_\text{v,n}$},
				ytick=\empty,
                xmin=0, xmax=54,
                ymin=0, ymax=1.2,
                grid style={line width=.1pt, draw=gray!10},
                legend style={at={(0.98,0.98)},anchor=north east, font=\scriptsize},
                legend columns=1,
            ]
                \addplot[ybar, bar width=1.523, fill opacity=0.5, draw=none, area legend, fill=red] table[x=bin_center,y expr=\thisrow{count_per_um2}/1468932,col sep=comma] {figures/methods/arbib-c3s-pores/rt_CLD_volume.csv};
                \addlegendentry{\acs*{CDF}};

                \addplot[ybar, bar width=1.05, fill opacity=0.5, draw=none, area legend, fill=blue] table[x=bin_center,y expr=\thisrow{count_per_um2}/43273,col sep=comma] {figures/methods/arbib-c3s-pores/rt_PSD_volume.csv};
                \addlegendentry{\acs*{PSD}};

				%\addplot [no markers, orange, domain=0:55, samples=100] {1142970 / 1468932 * exp(-0.5 * ((x - 17.4412) / 9.1613)^2)};

				%\addplot [no markers, cyan, domain=0:55, samples=100] {31149 / 43273 * exp(-0.5 * ((x - 17.8718) / 9.2256)^2)};

				\addplot [no markers, red, domain=0.01:55, samples=100] {(20042957.1493 / x * exp(-0.5 * ((ln(x) - 2.9845) / 0.5814)^2)) / 1468932};
				\addlegendentry{\acs*{CDF} fit}
				% max value at x=14.1038

				\addplot [no markers, blue, domain=0.01:55, samples=100] {(564085.0450 / x * exp(-0.5 * ((ln(x) - 3.0072) / 0.5615)^2)) / 43273};
				\addlegendentry{\acs*{PSD} fit}
				% max value at x=14.7599

            \end{axis}
        \end{tikzpicture}
        \caption{Room temperature (\qty{20}{\celsius}).}
        \label{fig:compare_rt_c}
    \end{subfigure}
    \hfill
    % Cryo Graph
    \begin{subfigure}{0.49\textwidth}
        \flushleft
        \begin{tikzpicture}
            \begin{axis}[
                width=1.1\linewidth,
                height=6cm,
                grid=major,
                xlabel={pore radius in \unit{\nano\meter}},
                ylabel={weighted frequency $f_\text{v,n}$},
                yticklabel pos=right,
				ytick=\empty,
                xmin=0, xmax=54,
                ymin=0, ymax=1.2,
                grid style={line width=.1pt, draw=gray!10},
                legend style={at={(0.98,0.98)},anchor=north east, font=\scriptsize},
            ]
                \addplot[ybar, bar width=1.023, fill opacity=0.5, draw=none, area legend, fill=red] table[x=bin_center,y expr=\thisrow{count_per_um2}/470539,col sep=comma] {figures/methods/arbib-c3s-pores/cryo_CLD_volume.csv};
                \addlegendentry{\acs*{CDF}};

                \addplot[ybar, bar width=0.574, fill opacity=0.5, draw=none, area legend, fill=blue] table[x=bin_center,y expr=\thisrow{count_per_um2}/9464,col sep=comma] {figures/methods/arbib-c3s-pores/cryo_PSD_volume.csv};
                \addlegendentry{\acs*{PSD}};

				%\addplot [no markers, orange, domain=0:40, samples=100] {386465/470539 * exp(-0.5 * ((x - 10.6056) / 5.0324)^2)};

				%\addplot [no markers, cyan, domain=0:40, samples=100] {7179/9464 * exp(-0.5 * ((x - 10.5568) / 4.5719)^2)};

				\addplot [no markers, red, domain=0.01:55, samples=100] {(4085894.9164 / x * exp(-0.5 * ((ln(x) - 2.4277) / 0.4976)^2))/470539.9};
				\addlegendentry{\acs*{CDF} fit}
				% max value at x=8.8474

				\addplot [no markers, blue, domain=0.01:55, samples=100] {(76152.7696 / x * exp(-0.5 * ((ln(x) - 2.4129) / 0.4371)^2))/9464};
				\addlegendentry{\acs*{PSD} fit}
				% max value at x=9.2243

            \end{axis}
        \end{tikzpicture}
        \caption{Cryo conditions (\qty{-140}{\celsius}).}
        \label{fig:compare_cryo_c}
    \end{subfigure}
	\caption{
		\ac{PSD} and \ac{CDF} of \ac{CSHi} (\qty{28}{\day} hydrated \ac{C3S}) prepared at \qty{20}{\celsius} and \qty{-140}{\celsius}, obtained from \ac{SE} images (\qty{20}{\celsius}: analysed \qty{23.2}{\micro\meter\squared}; \qty{-140}{\celsius}: analysed \qty{24.0}{\micro\meter\squared}).
		The fit parameters are listed in \zcref{tab:rP_fit_parameters}.
	}
    \label{fig:compare_graphs}
\end{figure}

As illustrated in \zcref{fig:compare_cryo_c,fig:compare_rt_c}, the \ac{CDF} (converted from chord to volume, red) exhibits a distribution remarkably similar to the \ac{PSD} determined by pore area analysis (blue).
However, the data density of the \ac{PSD} is relatively low, resulting in significant fluctuations.
The data density of the \ac{PSD} is also insufficient for evaluating individual images, as previously demonstrated in \cite{kleiner.2021a}.
Furthermore, the \ac{PSD} appears to be more sensitive to image background noise (compare \zcref{fig:compare_cryo_c,fig:compare_rt_c}), whereas the \ac{CDF} yields more robust results due to its higher data density.
The analysis indicates that the pores within \ac{CSHi} predominantly possess diameters in the range of \qtyrange{5}{30}{\nano\meter}.

The position of the most dominant pore radius, $r_\text{P, max}$, determined by the fit function (\zcref{eq:lognormal_fit}) of the \ac{PSD} in the range $r_\text{P}=\qtyrange{0}{55}{\nano\meter}$, was utilised to compare both methods.
The fit parameters are listed in \zcref{tab:rP_fit_parameters}.
While the individual evaluation of the images is not enough to reveal a significant influence of the temperature used for the \ac{ArBIB} milling \cite{kleiner.2021a}, the summarised dataset indicates a difference of the preparation temperature.
Combining the datasets, the pore radius was determined to be \qty{9.2}{\nano\meter} for \ac{ArBIB} at \qty{-140}{\celsius} and \qty{14.8}{\nano\meter} for \ac{ArBIB} at \qty{20}{\celsius}.
Furthermore, a larger pore area fraction was measured at room temperature compared to cryo conditions (\qty{4.1}{\percent} vs. \qty{1.3}{\percent}).
While a trend towards larger nano-porosity in the \ac{CSHi} prepared at room temperature is measurable, it must be noted that the analysed area was relatively small, and only a small portion contained segmented porosity.

Finally, \zcref{fig:compare_cryo_a,fig:compare_rt_a} reveal no apparent differences in the pore structure.
However, it appears that the curtaining artefacts visible in some images may affect the measured pore sizes.
%The observed variation in the \ac{CDF} between the two methods is negligible and likely attributable to variations in the positioning and quality of the acquired images.



\subsubsection{\acs*{FIBnt}}

To obtain data on the nano porosity of \ac{CSHi} of \qty{28}{\day} hydrated \ac{C3S}, a high-resolution \ac{BSE}-stack was created using \ac{FIBnt} (see \zcref{fig:fib_highres_pores_a}) \cite{kleiner.2026}.

The dataset has a volume of \qtyproduct{3.8 x 4.1 x 0.7}{\micro\meter\cubed} and a voxel size \qtyproduct{1.6 x 2.0 x 10}{\nano\meter\cubed}.
It contains a block of \ac{CSHi} which is surrounded by \ac{CSH} needles and includes a \textsc{Hadley} grain with an unhydrated \ac{C3S} particle in its centre.
\zcref{fig:fib_highres_pores_b} illustrates the porosity within the \ac{CSHi}.

% E:\Mitarbeiter\Florian Kleiner\FIB-Data\C3S 1Tag-10Monate\2019_11_01 C3S High Resolution Training RT
\begin{figure}[htb]
    \centering

	\begin{subfigure}[t]{0.43\textwidth}
		\ifdefined\PRINTABLE
			\includegraphics[width=\textwidth]{fibnt-datasets/high-res-ani/high-res-11.pdf}
		\else
			\animategraphics[palindrome,autoplay,poster=11,width=\textwidth]{10}{fibnt-datasets/high-res-ani/high-res-}{00}{23}
		\fi
		\caption{\ac{BSE} slice of the stack.}
		\label{fig:fib_highres_pores_a}
	\end{subfigure}
	\hfill
	\begin{subfigure}[t]{0.52\textwidth}
		\ifdefined\PRINTABLE
			\includegraphics[width=\textwidth]{fibnt-datasets/vol_ani/volume_00.pdf}
		\else
			\animategraphics[palindrome,autoplay,poster=00,width=\textwidth]{10}{fibnt-datasets/vol_ani/volume_}{00}{14}
		\fi
		\caption{3D-reconstruction of the pore space.}
		\label{fig:fib_highres_pores_b}
	\end{subfigure}
	\caption{
		\ac{FIBnt} dataset (\qtyproduct{3.8 x 4.1 x 0.7}{\micro\meter\cubed}) of a \qty{28}{\day} hydrated \ac{C3S} specimen and its 3D-reconstruction \cite{kleiner.2026}.
		Porosity (solid, grey) within \ac{CSHi} (semi-transparent, {\color{red}red}).
	}
	\label{fig:fib_highres_pores}
\end{figure}


While the energy density introduced in the specimen surface by imaging at \qty{2}{\kilo\volt} and \qty{0.1}{\nano\ampere} was fairly low, it is still possible that the nano-porosity observed within the bulk was induced or at least enlarged during the image acquisition as discussed in \zcref{sec:fib_evaluation_obstacles}.

Each slice of the volume was denoised using \ac{NLM} and segmented using a global threshold value.
The segmented pores were then evaluated.
Assuming spherical porosity, the pore radius $r_\text{P}$ was calculated for the segmented volume.
The respective pore size distribution histograms are shown in \zcref{fig:fib_highres_hist_3D}.

%X:\Mitarbeiter\Florian Kleiner\FIB-Data\C3S 1Tag-10Monate\2019_11_01 C3S High Resolution Training RT\Export-porosity\
\begin{figure}[H]
	\centering
	\begin{subfigure}[t]{0.49\textwidth}
		\raggedright
		\begin{tikzpicture}
			\begin{axis}[
				width=1.1\textwidth,
				height=6cm,
				xlabel={pore radius $r_\text{P}$ in \unit{\nano\meter}},
				ylabel={normalised frequency},
				grid=major,
				xmin=0, xmax=95,
				xtick={0,20,...,100},
				ymin=0,
				ytick=\empty,
				grid style={line width=.1pt, draw=gray!10},
				legend style={at={(0.98,0.98)}, anchor=north east},
				%axis on top,
			]
			% \addplot [ybar, bar width=1.4492, bar shift=0pt, fill=blue, fill opacity=.5, draw=none, area legend] table [x=Radius_Center, y expr=\thisrow{Frequency}/10929, col sep=comma] {figures/fibnt-datasets/pores-high-res_radius_histogram_2D.csv};
			% \addlegendentry{2D}
			\addplot [ybar, bar width=1.3795, bar shift=0pt, fill=blue, fill opacity=.5, draw=none, area legend] table [x=Radius_Center, y expr=\thisrow{Frequency}/2964, col sep=comma] {figures/fibnt-datasets/pores-high-res_radius_histogram_3D.csv};
			\addlegendentry{\ac{CSHi} porosity}
			\end{axis}
		\end{tikzpicture}
		\caption{\ac{PSD} of the 2D slices within \ac{CSHi}.}
		\label{fig:fib_highres_hist_3D-a}
	\end{subfigure}
	\hfill
	\begin{subfigure}[t]{0.49\textwidth}
		\raggedleft
		\begin{tikzpicture}
			\begin{axis}[
				width=1.1\textwidth,
				height=6cm,
				xlabel={pore radius $r_\text{P}$ in \unit{\nano\meter}},
				ylabel={normalised frequency},
				grid=major,
				xmin=0, xmax=95,
				xticklabels={0,20,...,100},
				xtick={0,20,...,100},
				ymin=0,ymax=1.2,
				ytick=\empty,
				grid style={line width=.1pt, draw=gray!10},
				yticklabel pos=right,
				legend style={at={(0.98,0.98)}, anchor=north east},
				%axis on top,
			]
			% \addplot [ybar, bar width=1.4492, bar shift=0pt, fill=blue, fill opacity=.5, draw=none, area legend] table [x=Radius_Center, y expr=\thisrow{Total_per_bin}/5011792.00, col sep=comma] {figures/fibnt-datasets/pores-high-res_bin_volumes_2D.csv};
			% \addlegendentry{2D}
			\addplot [ybar, bar width=1.3795, bar shift=0pt, fill=red, fill opacity=.5, draw=none, area legend, restrict x to domain=7:95] table [x=Radius_Center, y expr=\thisrow{Total_per_bin} / 1396002.62, col sep=comma] {figures/fibnt-datasets/pores-high-res_bin_volumes_3D.csv};
			\addlegendentry{volume porosity}

			% the normalisation of the fit is not correct for some reason
			%\addplot [no markers, red, domain=0.01:95, samples=100] {(23763476.64 / x * exp(-0.5 * ((ln(x) - 2.6867) / 0.5548)^2)) / 1396002.62};
			\addplot [no markers, red, domain=6.89:95, samples=100] { ((17087615.93 / x) * exp(-0.5 * ((ln(x) - 2.6867) / 0.5548)^2)) / 1357352.69 };
			\addlegendentry{fit}

			\addplot [ybar, bar width=1.3795, bar shift=0pt, fill=black, fill opacity=.3, draw=none, area legend, restrict x to domain=0:7] table [x=Radius_Center, y expr=\thisrow{Total_per_bin} / 1396002.62, col sep=comma] {figures/fibnt-datasets/pores-high-res_bin_volumes_3D.csv};

			\end{axis}
		\end{tikzpicture}
		\caption{\ac{PSD} of the volume within \ac{CSHi}.}
		\label{fig:fib_highres_hist_3D-b}
	\end{subfigure}
	\caption{
		Normalised \ac{PSD} of the pores measured using the segmented 2D slices ({\color{blue}blue}) and the \ac{PSD} of the segmented dataset based on the segmented 3D volume of of \ac{CSHi} ({\color{red}red}).
		% and the 2D slices ({\color{red}red}).
		The first three bins ({\color{darkgrey}grey}) were not used to for the fit.
	}
	\label{fig:fib_highres_hist_3D}
\end{figure}


The weighted and normalised frequency $f_\text{v,n}$ (see \zcref{eq:weig_norm_pore_frequency}) of pores increases with decreasing pore radius, with the smallest detectable pores being the most abundant.
However, the distribution differs when analysing the contribution to the total pore volume.
The peak of the pore volume distribution was determined using a log-normal fit function (\zcref{eq:lognormal_fit}) and was found at a pore radius $r_\text{P}$ of \qty{10.8}{\nano\meter}.
The fit parameters are listed in \zcref{tab:rP_fit_parameters}.
For this fit, only radii larger than \qty{6.5}{\nano\meter} were considered, as the three excluded smaller bins (marked in grey) are likely attributable to data noise.
Notably, the voxel geometry exhibits significant anisotropy, with a slice thickness ($z$) of \qty{10}{\nano\meter} compared to a lateral resolution ($x, y$) of \qtyproduct{1.6 x 2.0}{\nano\meter}.
This anisotropy introduces considerable uncertainty in measuring pore diameters below \qty{10}{\nano\meter}, as these features cannot be reliably resolved in the $z$-direction.

Considering these factors and the potential for pore enlargement caused by the preparation method, the measured distribution is likely shifted towards larger pore radii.
Consequently, the most frequent pore radius of the distribution is expected to be lower than the determined \qty{10.8}{\nano\meter}.

Comparing this value with the results from the \ac{ArBIB}-prepared samples in the previous section, it falls between the values measured for the two temperature conditions: \qty{9.2}{\nano\meter} at \qty{-140}{\celsius} and \qty{14.8}{\nano\meter} at room temperature.

% Consequently, a purely volumetric approach for such an anisotropic dataset presents challenges.
% Evaluating the dataset based on the segmented 2D slices yields a maximum pore volume at a radius of \qty{10.7}{\nano\meter}.
% However, the histogram extrapolated to volume data from 2D slice data also demonstrates, that a direct comparison is not directly possible (see \zcref{fig:fib_highres_hist_2D-b})
% Despite this shift in the peak value, the overall distribution of pore radii remains comparable between the two methods.



\FloatBarrier
\subsubsection{\acs{LTDSC}}

The \ac{LTDSC} measurements of the second freezing pass of \ac{C3S} and \ac{C2S} from \qtyrange{1}{28}{\day} of hydration are shown in \zcref{fig:ltdsc_result_freeze} (upper graphs \ac{C3S}, lower graphs \ac{C2S}) and available in \cite{kleiner.2026h}.
The graph shows the mean curve from three individual measurements for each specimen age.
The individual measurements can be found in \zcref{fig:ltdsc_result_freeze_raw}.
It was not possible to measure \ac{C2S} after \qty{1}{\day} of hydration, since the material was not hardened enough to be handled and transferred into the \ac{LTDSC} crucible.

\begin{figure}[!htb]
    \centering

 	\begin{subfigure}{.47\textwidth}
      \begin{tikzpicture}[remember picture]
        \begin{semilogxaxis}[
				blue,
				axis line style={black},
				xmin=1.5, xmax=60,
				ymin=-.20, ymax=0,
				y dir=reverse,
				width=\linewidth,
				height=6cm,
				grid=none,
				xlabel={temperature in \unit{\celsius}},
				xticklabels={\num{-51.2}, \num{-25.3}, \num{-12.4}, \num{-4.6}, \num{-2.0}},
				xtick={2.5,5,10,25,50},
				ytick=\empty,
				scaled y ticks=false,
				scaled x ticks=false,
				axis x line*=top,
            ]

        \end{semilogxaxis}
        \begin{semilogxaxis}[
				legend style={at={(0.95,.7)},anchor=east},
                legend cell align={left},
				xmin=1.5, xmax=60,
				ymin=-.20, ymax=0.07,
				y dir=reverse,
				width=\linewidth,
				height=6cm,
				grid=major,
				xlabel={pore radius $r_\text{P}$ \unit{\nano\meter}},
				ylabel={heat flow $\dot{Q}$ in \unit{\joule\per\second\per\gram}},
				grid style={line width=.1pt, draw=gray!10},
				yticklabel style={
						/pgf/number format/fixed,
						/pgf/number format/precision=1,
						/pgf/number format/fixed zerofill,
				},
				xticklabels={2.5,5,10,25,50},
				xtick={2.5,5,10,25,50},
				scaled y ticks=false,
				scaled x ticks=false,
				clip=false,
            ]
			\addplot[dotted, thick, samples=2, darkgrey,forget plot] coordinates {(3.19,-.2)(3.19,0.07)};
			\node[darkgrey, rotate=90, anchor=west] at (axis cs: 2.51,-.105) {\scriptsize <\,\qty{3.19}{\nano\meter}};

			% std dev of 1d C3S
            \addplot+ [draw=none, no markers, name path=A,forget plot] table [x index=3, col sep=comma,y expr={\thisrowno{1}-\thisrowno{2}}] {figures/ltdsc/reduced/C3S_1d_mean_freeze_segment.csv};
            \addplot+ [draw=none, no markers, name path=B,forget plot] table [x index=3, col sep=comma,y expr={\thisrowno{1}+\thisrowno{2}}] {figures/ltdsc/reduced/C3S_1d_mean_freeze_segment.csv};
            \addplot [fill=yellow, fill opacity=0.2,forget plot] fill between [of=A and B];

			% std dev of 7d C3S
            \addplot+ [draw=none, no markers, name path=A,forget plot] table [x index=3, col sep=comma,y expr={\thisrowno{1}-\thisrowno{2}}] {figures/ltdsc/reduced/C3S_7d_mean_freeze_segment.csv};
            \addplot+ [draw=none, no markers, name path=B,forget plot] table [x index=3, col sep=comma,y expr={\thisrowno{1}+\thisrowno{2}}] {figures/ltdsc/reduced/C3S_7d_mean_freeze_segment.csv};
            \addplot [fill=orange, fill opacity=0.2,forget plot] fill between [of=A and B];

			% std dev of 14d C3S
            \addplot+ [draw=none, no markers, name path=A,forget plot] table [x index=3, col sep=comma,y expr={\thisrowno{1}-\thisrowno{2}}] {figures/ltdsc/reduced/C3S_14d_mean_freeze_segment.csv};
            \addplot+ [draw=none, no markers, name path=B,forget plot] table [x index=3, col sep=comma,y expr={\thisrowno{1}+\thisrowno{2}}] {figures/ltdsc/reduced/C3S_14d_mean_freeze_segment.csv};
            \addplot [fill=red, fill opacity=0.2,forget plot] fill between [of=A and B];

			% std dev of 21d C3S
            \addplot+ [draw=none, no markers, name path=A,forget plot] table [x index=3, col sep=comma,y expr={\thisrowno{1}-\thisrowno{2}}] {figures/ltdsc/reduced/C3S_21d_mean_freeze_segment.csv};
            \addplot+ [draw=none, no markers, name path=B,forget plot] table [x index=3, col sep=comma,y expr={\thisrowno{1}+\thisrowno{2}}] {figures/ltdsc/reduced/C3S_21d_mean_freeze_segment.csv};
            \addplot [fill=green, fill opacity=0.2,forget plot] fill between [of=A and B];

			% std dev of 28d C3S
            \addplot+ [draw=none, no markers, name path=A,forget plot] table [x index=3, col sep=comma,y expr={\thisrowno{1}-\thisrowno{2}}] {figures/ltdsc/reduced/C3S_28d_mean_freeze_segment.csv};
            \addplot+ [draw=none, no markers, name path=B,forget plot] table [x index=3, col sep=comma,y expr={\thisrowno{1}+\thisrowno{2}}] {figures/ltdsc/reduced/C3S_28d_mean_freeze_segment.csv};
            \addplot [fill=blue, fill opacity=0.2,forget plot] fill between [of=A and B];


            \addplot [ no marks, yellow] table [x index=3, y index=1, col sep=comma] {figures/ltdsc/reduced/C3S_1d_mean_freeze_segment.csv};
            %\addlegendentry{\ac{C3S}, \qty{1}{\day}};
            \addplot [ no marks, orange] table [x index=3, y index=1, col sep=comma] {figures/ltdsc/reduced/C3S_7d_mean_freeze_segment.csv};
            %\addlegendentry{\ac{C3S}, \qty{7}{\day}};
            \addplot [ no marks, red] table [x index=3, y index=1, col sep=comma] {figures/ltdsc/reduced/C3S_14d_mean_freeze_segment.csv};
            %\addlegendentry{\ac{C3S}, \qty{14}{\day}};
            \addplot [ no marks, green] table [x index=3, y index=1, col sep=comma] {figures/ltdsc/reduced/C3S_21d_mean_freeze_segment.csv};
            %\addlegendentry{\ac{C3S}, \qty{21}{\day}};
            \addplot [ no marks, blue] table [x index=3, y index=1, col sep=comma] {figures/ltdsc/reduced/C3S_28d_mean_freeze_segment.csv};
            %\addlegendentry{\ac{C3S}, \qty{28}{\day}};

			\draw[draw=black, densely dotted] (2.75,-.11) rectangle (1.75,-0.03);
			\node[draw=none] (top-3) at (axis cs:2.75,-.11) {};
			\node[draw=none] (bottom-3) at (axis cs:2.75,-.03) {};

            \node[anchor=west, align=left] at (.8,-0.25) {\textbf{\ac{C3S}}};
        \end{semilogxaxis}
	  \end{tikzpicture}

	\end{subfigure}
	\hfill
	\begin{subfigure}{.47\textwidth}

	  \begin{tikzpicture}[remember picture]
        \begin{axis}[
			blue,
			axis line style={black},
			xmin=1.75, xmax=2.75,
			ymin=-.11, ymax=-0.03,
			y dir=reverse,
			width=.85\linewidth,
			height=6cm,
			grid=none,
			xlabel={temperature in \unit{\celsius}},
			xticklabels={\num{-73.3}, \num{-64.1}, \num{-56.9}, \num{-51.2}, \num{-46.5}},
			xtick={1.75,2.00,2.25,2.50,2.75},
			ytick=\empty,
			scaled y ticks=false,
			scaled x ticks=false,,
			axis x line*=top
		]

		\end{axis}
		\begin{axis}[
				legend style={at={(0.97,.97)},anchor=north east},
				legend cell align={left},
				xmin=1.75, xmax=2.75,
				ymin=-.11, ymax=-0.03,
				y dir=reverse,
				width=.85\linewidth,
				height=6cm,
				grid=major,
				xlabel={pore radius $r_\text{P}$ \unit{\nano\meter}},
				ylabel={heat flow $\dot{Q}$ in \unit{\joule\per\second\per\gram}},
				grid style={line width=.1pt, draw=gray!10},
				xticklabels={1.75,2.00,2.25,2.50,2.75},
				xtick={1.75,2.00,2.25,2.50,2.75},
				yticklabel style={
						/pgf/number format/fixed,
						/pgf/number format/precision=2,
						/pgf/number format/fixed zerofill,
				},
				scaled y ticks=false,
				scaled x ticks=false,
				yticklabel pos=right,
			]

			% std dev of 1d C3S
            \addplot+ [draw=none, no markers, name path=A,forget plot] table [x index=3, col sep=comma,y expr={\thisrowno{1}-\thisrowno{2}}] {figures/ltdsc/reduced/zoom/C3S_1d_mean_freeze_segment.csv};
            \addplot+ [draw=none, no markers, name path=B,forget plot] table [x index=3, col sep=comma,y expr={\thisrowno{1}+\thisrowno{2}}] {figures/ltdsc/reduced/zoom/C3S_1d_mean_freeze_segment.csv};
            \addplot [fill=yellow, fill opacity=0.2,forget plot] fill between [of=A and B];

			% std dev of 7d C3S
            \addplot+ [draw=none, no markers, name path=A,forget plot] table [x index=3, col sep=comma,y expr={\thisrowno{1}-\thisrowno{2}}] {figures/ltdsc/reduced/zoom/C3S_7d_mean_freeze_segment.csv};
            \addplot+ [draw=none, no markers, name path=B,forget plot] table [x index=3, col sep=comma,y expr={\thisrowno{1}+\thisrowno{2}}] {figures/ltdsc/reduced/zoom/C3S_7d_mean_freeze_segment.csv};
            \addplot [fill=orange, fill opacity=0.2,forget plot] fill between [of=A and B];

			% std dev of 14d C3S
            \addplot+ [draw=none, no markers, name path=A,forget plot] table [x index=3, col sep=comma,y expr={\thisrowno{1}-\thisrowno{2}}] {figures/ltdsc/reduced/zoom/C3S_14d_mean_freeze_segment.csv};
            \addplot+ [draw=none, no markers, name path=B,forget plot] table [x index=3, col sep=comma,y expr={\thisrowno{1}+\thisrowno{2}}] {figures/ltdsc/reduced/zoom/C3S_14d_mean_freeze_segment.csv};
            \addplot [fill=red, fill opacity=0.2,forget plot] fill between [of=A and B];

			% std dev of 21d C3S
            \addplot+ [draw=none, no markers, name path=A,forget plot] table [x index=3, col sep=comma,y expr={\thisrowno{1}-\thisrowno{2}}] {figures/ltdsc/reduced/zoom/C3S_21d_mean_freeze_segment.csv};
            \addplot+ [draw=none, no markers, name path=B,forget plot] table [x index=3, col sep=comma,y expr={\thisrowno{1}+\thisrowno{2}}] {figures/ltdsc/reduced/zoom/C3S_21d_mean_freeze_segment.csv};
            \addplot [fill=green, fill opacity=0.2,forget plot] fill between [of=A and B];

			% std dev of 28d C3S
            \addplot+ [draw=none, no markers, name path=A,forget plot] table [x index=3, col sep=comma,y expr={\thisrowno{1}-\thisrowno{2}}] {figures/ltdsc/reduced/zoom/C3S_28d_mean_freeze_segment.csv};
            \addplot+ [draw=none, no markers, name path=B,forget plot] table [x index=3, col sep=comma,y expr={\thisrowno{1}+\thisrowno{2}}] {figures/ltdsc/reduced/zoom/C3S_28d_mean_freeze_segment.csv};
            \addplot [fill=blue, fill opacity=0.2,forget plot] fill between [of=A and B];

			% mean values of the 1d C3S
			\addplot+[no marks, orange] table [x index=3, y index=1, col sep=comma] {figures/ltdsc/reduced/zoom/C3S_1d_mean_freeze_segment.csv};
			\addlegendentry{\qty{1}{\day}};
			% mean values of the 7d C3S
			\addplot+[no marks, orange] table [x index=3, y index=1, col sep=comma] {figures/ltdsc/reduced/zoom/C3S_7d_mean_freeze_segment.csv};
			\addlegendentry{\qty{7}{\day}};
			% mean values of the 14d C3S
			\addplot+[no marks, red] table [x index=3, y index=1, col sep=comma] {figures/ltdsc/reduced/zoom/C3S_14d_mean_freeze_segment.csv};
			\addlegendentry{\qty{14}{\day}};
			% mean values of the 21d C3S
			\addplot+[no marks, green] table [x index=3, y index=1, col sep=comma] {figures/ltdsc/reduced/zoom/C3S_21d_mean_freeze_segment.csv};
			\addlegendentry{\qty{21}{\day}};
			% mean values of the 28d C3S
			\addplot+[no marks, blue] table [x index=3, y index=1, col sep=comma] {figures/ltdsc/reduced/zoom/C3S_28d_mean_freeze_segment.csv};
			\addlegendentry{\qty{28}{\day}};

            %\node[anchor=west, align=left] at (axis cs:2,-0.09) {\ac{C3S}};

			\node[draw=none] (top-4) at (axis cs:1.75,-.11) {};
			\node[draw=none] (bottom-4) at (axis cs:1.75,-.03) {};
		\end{axis}
	  \end{tikzpicture}
	\end{subfigure}
    \begin{tikzpicture}[remember picture,overlay]
        \draw[dashed,grey] (top-3) -- (top-4);
        \draw[dashed,grey] (bottom-3) -- (bottom-4);
    \end{tikzpicture}

	\vspace{1\baselineskip}

 	\begin{subfigure}{.47\textwidth}
      \begin{tikzpicture}[remember picture]
        \begin{semilogxaxis}[
				blue,
				axis line style={black},
				xmin=1.5, xmax=60,
				ymin=-.20, ymax=0,
				y dir=reverse,
				width=\linewidth,
				height=6cm,
				grid=none,
				xlabel={temperature in \unit{\celsius}},
				xticklabels={\num{-51.2}, \num{-25.3}, \num{-12.4}, \num{-4.6}, \num{-2.0}},
				xtick={2.5,5,10,25,50},
				ytick=\empty,
				scaled y ticks=false,
				scaled x ticks=false,,
				axis x line*=top
            ]

        \end{semilogxaxis}
        \begin{semilogxaxis}[
				legend style={at={(0.95,.7)},anchor=east},
                legend cell align={left},
				xmin=1.5, xmax=60,
				ymin=-.20, ymax=0.07,
				y dir=reverse,
				width=\linewidth,
				height=6cm,
				grid=major,
				xlabel={pore radius $r_\text{P}$ \unit{\nano\meter}},
				ylabel={heat flow $\dot{Q}$ in \unit{\joule\per\second\per\gram}},
				grid style={line width=.1pt, draw=gray!10},
				yticklabel style={
						/pgf/number format/fixed,
						/pgf/number format/fixed zerofill,
						/pgf/number format/precision=1
				},
				xticklabels={2.5,5,10,25,50},
				xtick={2.5,5,10,25,50},
				scaled y ticks=false,
				scaled x ticks=false,
				clip=false,
            ]
			\addplot[dotted, thick, samples=2, darkgrey,forget plot] coordinates {(3.19,-.2)(3.19,0.07)};
			\node[darkgrey, rotate=90, anchor=west] at (axis cs: 2.51,-.105) {\scriptsize <\,\qty{3.19}{\nano\meter}};

			% std dev of 7d C2S
            \addplot+ [draw=none, no markers, name path=A,forget plot] table [x index=3, col sep=comma,y expr={\thisrowno{1}-\thisrowno{2}}] {figures/ltdsc/reduced/C2S_7d_mean_freeze_segment.csv};
            \addplot+ [draw=none, no markers, name path=B,forget plot] table [x index=3, col sep=comma,y expr={\thisrowno{1}+\thisrowno{2}}] {figures/ltdsc/reduced/C2S_7d_mean_freeze_segment.csv};
            \addplot [fill=orange, fill opacity=0.2,forget plot] fill between [of=A and B];

			% std dev of 14d C2S
            \addplot+ [draw=none, no markers, name path=A,forget plot] table [x index=3, col sep=comma,y expr={\thisrowno{1}-\thisrowno{2}}] {figures/ltdsc/reduced/C2S_14d_mean_freeze_segment.csv};
            \addplot+ [draw=none, no markers, name path=B,forget plot] table [x index=3, col sep=comma,y expr={\thisrowno{1}+\thisrowno{2}}] {figures/ltdsc/reduced/C2S_14d_mean_freeze_segment.csv};
            \addplot [fill=red, fill opacity=0.2,forget plot] fill between [of=A and B];

			% std dev of 21d C2S
            \addplot+ [draw=none, no markers, name path=A,forget plot] table [x index=3, col sep=comma,y expr={\thisrowno{1}-\thisrowno{2}}] {figures/ltdsc/reduced/C2S_21d_mean_freeze_segment.csv};
            \addplot+ [draw=none, no markers, name path=B,forget plot] table [x index=3, col sep=comma,y expr={\thisrowno{1}+\thisrowno{2}}] {figures/ltdsc/reduced/C2S_21d_mean_freeze_segment.csv};
            \addplot [fill=green, fill opacity=0.2,forget plot] fill between [of=A and B];

			% std dev of 28d C2S
            \addplot+ [draw=none, no markers, name path=A,forget plot] table [x index=3, col sep=comma,y expr={\thisrowno{1}-\thisrowno{2}}] {figures/ltdsc/reduced/C2S_28d_mean_freeze_segment.csv};
            \addplot+ [draw=none, no markers, name path=B,forget plot] table [x index=3, col sep=comma,y expr={\thisrowno{1}+\thisrowno{2}}] {figures/ltdsc/reduced/C2S_28d_mean_freeze_segment.csv};
            \addplot [fill=blue, fill opacity=0.2,forget plot] fill between [of=A and B];


            \addplot [ no marks, orange] table [x index=3, y index=1, col sep=comma] {figures/ltdsc/reduced/C2S_7d_mean_freeze_segment.csv};
            %\addlegendentry{\ac{C2S}, \qty{7}{\day}};
            \addplot [ no marks, red] table [x index=3, y index=1, col sep=comma] {figures/ltdsc/reduced/C2S_14d_mean_freeze_segment.csv};
            %\addlegendentry{\ac{C2S}, \qty{14}{\day}};
            \addplot [ no marks, green] table [x index=3, y index=1, col sep=comma] {figures/ltdsc/reduced/C2S_21d_mean_freeze_segment.csv};
            %\addlegendentry{\ac{C2S}, \qty{21}{\day}};
            \addplot [ no marks, blue] table [x index=3, y index=1, col sep=comma] {figures/ltdsc/reduced/C2S_28d_mean_freeze_segment.csv};
            %\addlegendentry{\ac{C2S}, \qty{28}{\day}};

			\draw[draw=black, densely dotted] (2.75,-.11) rectangle (1.75,-0.03);
			\node[draw=none] (top-1) at (axis cs:2.75,-.11) {};
			\node[draw=none] (bottom-1) at (axis cs:2.75,-.03) {};

            \node[anchor=west, align=left] at (.8,-0.25) {\textbf{\ac{C2S}}};

        \end{semilogxaxis}
	  \end{tikzpicture}

	\end{subfigure}
	\hfill
	\begin{subfigure}{.47\textwidth}

	  \begin{tikzpicture}[remember picture]
        \begin{axis}[
			blue,
			axis line style={black},
			xmin=1.75, xmax=2.75,
			ymin=-.11, ymax=-0.03,
			y dir=reverse,
			width=.85\linewidth,
			height=6cm,
			grid=none,
			xlabel={temperature in \unit{\celsius}},
			xticklabels={\num{-73.3}, \num{-64.1}, \num{-56.9}, \num{-51.2}, \num{-46.5}},
			xtick={1.75,2.00,2.25,2.50,2.75},
			ytick=\empty,
			scaled y ticks=false,
			scaled x ticks=false,,
			axis x line*=top
		]

		\end{axis}
		\begin{axis}[
				legend style={at={(0.97,.97)},anchor=north east},
				legend cell align={left},
				xmin=1.75, xmax=2.75,
				ymin=-.11, ymax=-0.03,
				y dir=reverse,
				width=.85\linewidth,
				height=6cm,
				grid=major,
				xlabel={pore radius $r_\text{P}$ \unit{\nano\meter}},
				ylabel={heat flow $\dot{Q}$ in \unit{\joule\per\second\per\gram}},
				grid style={line width=.1pt, draw=gray!10},
				xticklabels={1.75,2.00,2.25,2.50,2.75},
				xtick={1.75,2.00,2.25,2.50,2.75},
				yticklabel style={
						/pgf/number format/fixed,
						/pgf/number format/precision=2,
						/pgf/number format/fixed zerofill,
				},
				scaled y ticks=false,
				scaled x ticks=false,
				yticklabel pos=right,
			]
			% std dev of 7d C2S
            \addplot+ [draw=none, no markers, name path=A,forget plot] table [x index=3, col sep=comma,y expr={\thisrowno{1}-\thisrowno{2}}] {figures/ltdsc/reduced/zoom/C2S_7d_mean_freeze_segment.csv};
            \addplot+ [draw=none, no markers, name path=B,forget plot] table [x index=3, col sep=comma,y expr={\thisrowno{1}+\thisrowno{2}}] {figures/ltdsc/reduced/zoom/C2S_7d_mean_freeze_segment.csv};
            \addplot [fill=orange, fill opacity=0.2,forget plot] fill between [of=A and B];

			% std dev of 14d C2S
            \addplot+ [draw=none, no markers, name path=A,forget plot] table [x index=3, col sep=comma,y expr={\thisrowno{1}-\thisrowno{2}}] {figures/ltdsc/reduced/zoom/C2S_14d_mean_freeze_segment.csv};
            \addplot+ [draw=none, no markers, name path=B,forget plot] table [x index=3, col sep=comma,y expr={\thisrowno{1}+\thisrowno{2}}] {figures/ltdsc/reduced/zoom/C2S_14d_mean_freeze_segment.csv};
            \addplot [fill=red, fill opacity=0.2,forget plot] fill between [of=A and B];

			% std dev of 21d C2S
            \addplot+ [draw=none, no markers, name path=A,forget plot] table [x index=3, col sep=comma,y expr={\thisrowno{1}-\thisrowno{2}}] {figures/ltdsc/reduced/zoom/C2S_21d_mean_freeze_segment.csv};
            \addplot+ [draw=none, no markers, name path=B,forget plot] table [x index=3, col sep=comma,y expr={\thisrowno{1}+\thisrowno{2}}] {figures/ltdsc/reduced/zoom/C2S_21d_mean_freeze_segment.csv};
            \addplot [fill=green, fill opacity=0.2,forget plot] fill between [of=A and B];

			% std dev of 28d C2S
            \addplot+ [draw=none, no markers, name path=A,forget plot] table [x index=3, col sep=comma,y expr={\thisrowno{1}-\thisrowno{2}}] {figures/ltdsc/reduced/zoom/C2S_28d_mean_freeze_segment.csv};
            \addplot+ [draw=none, no markers, name path=B,forget plot] table [x index=3, col sep=comma,y expr={\thisrowno{1}+\thisrowno{2}}] {figures/ltdsc/reduced/zoom/C2S_28d_mean_freeze_segment.csv};
            \addplot [fill=blue, fill opacity=0.2,forget plot] fill between [of=A and B];

			% mean values of the 7d C2S
			\addplot+[no marks, orange] table [each nth point=3, filter discard warning=false, unbounded coords=discard, x index=3, y index=1, col sep=comma] {figures/ltdsc/reduced/zoom/C2S_7d_mean_freeze_segment.csv};
			\addlegendentry{\qty{7}{\day}};
			% mean values of the 14d C2S
			\addplot+[no marks, red] table [each nth point=3, filter discard warning=false, unbounded coords=discard,x index=3, y index=1, col sep=comma] {figures/ltdsc/reduced/zoom/C2S_14d_mean_freeze_segment.csv};
			\addlegendentry{\qty{14}{\day}};
			% mean values of the 21d C2S
			\addplot+[no marks, green] table [each nth point=3, filter discard warning=false, unbounded coords=discard,x index=3, y index=1, col sep=comma] {figures/ltdsc/reduced/zoom/C2S_21d_mean_freeze_segment.csv};
			\addlegendentry{\qty{21}{\day}};
			% mean values of the 28d C2S
			\addplot+[no marks, blue] table [each nth point=3, filter discard warning=false, unbounded coords=discard,x index=3, y index=1, col sep=comma] {figures/ltdsc/reduced/zoom/C2S_28d_mean_freeze_segment.csv};
			\addlegendentry{\qty{28}{\day}};


			\node[draw=none] (top-2) at (axis cs:1.75,-.11) {};
			\node[draw=none] (bottom-2) at (axis cs:1.75,-.03) {};
		\end{axis}
	  \end{tikzpicture}
	\end{subfigure}
    \begin{tikzpicture}[remember picture,overlay]
        \draw[dashed,grey] (top-1) -- (top-2);
        \draw[dashed,grey] (bottom-1) -- (bottom-2);
    \end{tikzpicture}

    \caption{
		\ac{LTDSC} measurements (mean of three measurements) for the second freezing pass of \ac{C3S} (top) and \ac{C2S} (bottom) from \qtyrange{1}{28}{\day} of hydration \cite{kleiner.2026h}.
		Every specimen age was measured using three specimens.
		On the left-hand side, the full measurement is plotted (pore radius and temperature is plotted logarithmic).
		On the right-hand side is a magnified section (linear plot), highlighting the data correlating to the gel pores.
		Note that pore sizes below \qty{3.19}{\nano\meter} (<\,\qty{-40}{\celsius}) are undefined as noted in \zcref{eq:brun_freeze}.
	}
    \label{fig:ltdsc_result_freeze}
\end{figure}

The graphs on the right-hand side of \zcref{fig:ltdsc_result_freeze} are a magnified sections of the left graphs, showing the peaks correlating to the gel pores.
Applying \zcref{eq:brun_freeze}, a main pore radius range between \qtyrange{1.8}{2.3}{\nano\meter} can be calculated, while \ac{C3S} displays a slightly broader and larger pore size compared to \ac{C2S}.
However, according to \textcite{brun.1977} this equation is not valid for temperatures below \qty{-40}{\celsius}.
The figure shows that the heat flow signal of the gel pores is increasing with the age of the specimen, which indicates an increase in the gel pore volume with increasing specimen age.

At first glance, the \ac{C3S} results look similar to the \ac{C2S} results.
However, they have a lower signal response in the temperature range corresponding to the gel pores.
This indicates a lower gel pore volume for \ac{C3S} compared to \ac{C2S}.
Furthermore, the peaks are slightly shifted to smaller pore radii indicating smaller gel pores.
However, the freezing curve cannot be used to determine the volume of the gel pores.

In order to calculate the ice mass within the gel pores and thus also their volume, the heat flow signal of the two melting curves of every experiment was integrated over the measured temperature ranges and subtracted from one another, as shown in \zcref{eq:ltdsc_ice_integral, eq:ltdsc_ice2}.
Beforehand, a base line correction was applied to the raw data.
\zcref{fig:ltdsc_result_c2s_melt} illustrates this process for the results of \qty{28}{\day} hydrated \ac{C2S}.
%{\color{red}The detailed procedure is published on github \cite{kleiner.2026f} %\url{https://github.com/kleinerELM/Kalorimeterauswertung}
%}

Raw data are plotted as dashed lines, whereas the baseline-corrected data are shown as solid lines.
The corresponding baseline functions, $\dot{Q}_\text{B, P1}(T)$ and $\dot{Q}_\text{B, P2}(T)$, are also plotted in the graph.
Areas under the solid curves are highlighted in green and yellow.
Green shading represents the gel pore volume, while yellow indicates the capillary pore volume.
It should be noted that the green area fully overlaps the yellow area.

\begin{figure}[!htb]
    \centering

    \begin{tikzpicture}
        \begin{axis}[
				legend style={at={(0.1,.95)},anchor=north west},
                legend cell align={left},
				xmin=-60, xmax=15,
				ymin=0, ymax=0.6,
				width=\linewidth,
				height=7cm,
				grid=major,
				xlabel={temperature $\vartheta$ in \unit{\celsius}},
				ylabel={heat flow $\dot{Q}$ in \unit{\joule\per\second\per\gram}},
				xtick={-60, -50,...,10},
				grid style={line width=.1pt, draw=gray!10},
				yticklabel style={
						/pgf/number format/fixed,
						/pgf/number format/fixed zerofill,
						/pgf/number format/precision=1
				},
				scaled y ticks=false,
				scaled x ticks=false
            ]


            \addplot[fill=red, fill opacity=0.25, draw=none] coordinates {
                (-57, 0) (-56,0) (-56, 0.6) (-57,0.6) (-57,0)
            };
            \addplot[fill=red, fill opacity=0.25, draw=none] coordinates {
                (-21.5, 0) (-20.5,0) (-20.5, 0.6) (-21.5,0.6) (-21.5,0)
            };

            \addplot[fill=red, fill opacity=0.25, draw=none] coordinates {
                (9, 0) (10,0) (10, 0.6) (9,0.6) (9,0)
            };

            \node[anchor=west, align=left, font=\footnotesize]
                at (axis cs:-53,0.225) {
					$\dot{Q}_\text{B, P1}(T) = 0.00091T + 0.04939$\\
					$\dot{Q}_\text{B, P2}(T) = 0.00073T + 0.04064$};

            \node[anchor=west, align=left] at (axis cs:-4,0.5) {A};
            \node[anchor=west, align=left] at (axis cs:3.5,0.55) {B};

            \addplot [no marks, green, dashed] table [each nth point=4, filter discard warning=false, unbounded coords=discard,x index=0, y index=1, col sep=comma] {figures/ltdsc/ExpDat_C2S_28d-1_melt_segment_8.csv};
            \addlegendentry{P1, raw data};

            \addplot [no marks, blue, dashed] table [each nth point=4, filter discard warning=false, unbounded coords=discard,x index=0, y index=1, col sep=comma] {figures/ltdsc/ExpDat_C2S_28d-1_melt_segment_16.csv};
            \addlegendentry{P2, raw data};


			\addplot[thick, domain=0:1, color=blue, fill = yellow, fill opacity=0.3, forget plot] table [each nth point=4, filter discard warning=false, unbounded coords=discard,x index=0, y index=2, col sep=comma] {figures/ltdsc/ExpDat_C2S_28d-1_melt_segment_16.csv};

			\addplot[domain=0:1, color=green, fill = green, fill opacity=0.3, forget plot] table [each nth point=4, filter discard warning=false, unbounded coords=discard,x index=0, y index=2, col sep=comma] {figures/ltdsc/ExpDat_C2S_28d-1_melt_segment_8.csv};


            \addplot [thick, no marks, green] table [each nth point=4, filter discard warning=false, unbounded coords=discard,x index=0, y index=2, col sep=comma] {figures/ltdsc/ExpDat_C2S_28d-1_melt_segment_8.csv};
            \addlegendentry{P1, baseline corrected};

            \addplot [thick, no marks, blue] table [each nth point=4, filter discard warning=false, unbounded coords=discard,x index=0, y index=2, col sep=comma] {figures/ltdsc/ExpDat_C2S_28d-1_melt_segment_16.csv};
            \addlegendentry{P2, baseline corrected};

        \end{axis}
	\end{tikzpicture}

    \caption{
		\ac{LTDSC} curves for the first and second melting pass of \qty{28}{\day} hydrated \ac{C2S}.
		The {\color{green}green} area correlate to the capillary porosity, while the {\color{orange}yellow} area correlates to the gel porosity.
		The value range used for the linear baseline regression is marked {\color{red}red} between \qtyrange{-57}{-56}{\celsius} (P2) or \qtyrange{-21.5}{-20.5}{\celsius} (P1) and \qtyrange{9}{10}{\celsius} (P1 and P2).
	}
    \label{fig:ltdsc_result_c2s_melt}
\end{figure}


The melting curve shows two distinct peaks, one at \qty{0}{\celsius} (A) and one at \qty{2}{\celsius} (B).
This indicates that the heating rate of \qty{2}{\kelvin\per\minute} was higher than the inertia of the melting process, as the ice should melt at or before \qty{0}{\celsius}, but is delayed due to the fast temperature change (\qty{2}{\kelvin\per\minute}).

Therefore, an experiment for comparison with a slower heating rate (\qty{0.5}{\kelvin\per\minute}) was carried out (\zcref{fig:ltdsc_result_heating_rate_a}).
This resulted in a shift of the melting peak in the direction of \qty{0}{\celsius} and a seemingly smaller heat output.
However, the difference in peak height is caused by a change in the measuring interval.
To mitigate this, the dashed curve indicates the heat flow $\dot{Q}$ of the \qty{0.5}{\kelvin\per\minute} heat rate, corrected for the interval difference. %TODO: ist das legitim?

\begin{figure}[!htb]
    \centering
	\begin{subfigure}{.485\textwidth}
      \begin{tikzpicture}
        \begin{axis}[
			legend style={at={(0.05,.95)},anchor=north west,font=\footnotesize},
                legend cell align={left},
				xmin=-50, xmax=10,
				ymin=-0.01, ymax=0.2,
				width=\linewidth,
				height=7cm,
				grid=major,
				xlabel={temperature $\vartheta$ in \unit{\celsius}},
				ylabel={heat flow $\dot{Q}$ in \unit{\joule\per\second\per\gram}},
				grid style={line width=.1pt, draw=gray!10},
				yticklabel style={
						/pgf/number format/fixed,
						/pgf/number format/fixed zerofill,
						/pgf/number format/precision=2
				},
				scaled y ticks=false,
				scaled x ticks=false
            ]

			\addplot[samples=2, darkgrey,forget plot] coordinates {(15,0)(-60,0)};
			\addplot[samples=2, darkgrey,forget plot] coordinates {(0,-.01)(0,.2)};

            \addplot [thick, no marks, blue] table [each nth point=4, filter discard warning=false, unbounded coords=discard,x index=0, y index=2, col sep=comma] {figures/ltdsc/ExpDat_M28n2-28d_melt_segment_9.csv};
            \addlegendentry{\qty{0.5}{\kelvin\per\minute}};

			\addplot [dashed, no marks, blue] table [each nth point=4, filter discard warning=false, unbounded coords=discard, x index=0, y expr=\thisrowno{2}*4, col sep=comma] {figures/ltdsc/ExpDat_M28n2-28d_melt_segment_9.csv};
			\addlegendentry{\qty{0.5}{\kelvin\per\minute}$\,\times\,4$};


            \addplot [thick, no marks, red] table [each nth point=4, filter discard warning=false, unbounded coords=discard,x index=0, y index=2, col sep=comma] {figures/ltdsc/ExpDat_M28n2-28d_melt_segment_18.csv};
            \addlegendentry{\qty{2.0}{\kelvin\per\minute}};
        \end{axis}
	  \end{tikzpicture}
	  \caption{Melting heat flow}
	  \label{fig:ltdsc_result_heating_rate_a}

	\end{subfigure}
	\hfill
	\begin{subfigure}{.485\textwidth}

	  \begin{tikzpicture}[]
		\begin{axis}[
				legend style={at={(0.05,.95)},anchor=north west,font=\footnotesize},
				legend cell align={left},
				xmin=-50, xmax=10,
				ymin=-10, ymax=180,
				width=\linewidth,
				height=7cm,
				grid=major,
				xlabel={temperature $\vartheta$ in \unit{\celsius}},
				ylabel={cum. ice volume in \unit{\micro\liter\per\gram}},
				grid style={line width=.1pt, draw=gray!10},
				yticklabel style={
					/pgf/number format/fixed,
					/pgf/number format/precision=0
				},
				scaled y ticks=false,
				scaled x ticks=false,
				yticklabel pos=right,
			]
			\addplot[samples=2, darkgrey,forget plot] coordinates {(15,0)(-60,0)};

            \addplot [no marks, cyan] table [each nth point=10, filter discard warning=false, unbounded coords=discard,x index=0, y expr={\thisrowno{4}*1000}, col sep=comma] {figures/ltdsc/ExpDat_M28n2-28d_ice_9.csv};
            \addlegendentry{\qty{0.5}{\kelvin\per\minute}};

            \addplot [no marks, orange] table [each nth point=10, filter discard warning=false, unbounded coords=discard,x index=0,y expr={\thisrowno{4}*1000}, col sep=comma] {figures/ltdsc/ExpDat_M28n2-28d_ice_18.csv};
            \addlegendentry{\qty{2.0}{\kelvin\per\minute}};

		\end{axis}
	  \end{tikzpicture}
	  \caption{Cumulative ice volume}
	  \label{fig:ltdsc_result_heating_rate_b}

	\end{subfigure}

    \caption{
		\ac{LTDSC} curves and the cumulative ice volume per specimen weight for a slow (\qty{0.5}{\kelvin\per\minute}) and fast (\qty{2}{\kelvin\per\minute}) heating curves from \qtyrange{-50}{15}{\celsius} of \qty{28}{\day} hydrated cement (\ac{wb} = \num{0.5}).
		To better illustrate the effect, the heat flow signal of the \qty{0.5}{\kelvin\per\minute} measurement was also multiplied with a factor of \num{4} (dashed line).
	}
    \label{fig:ltdsc_result_heating_rate}
\end{figure}

% ,unbehau.2025 <- das Paper ist leider nicht rechtzeitig fertig geworden, müsste aber in \textcite{mueller.2021,unbehau.2025}...
The ice mass should not be affected by the heating rate.
The simplified approach to calculate the ice volume as used by \textcite{mueller.2021} (\zcref{eq:ltdsc_ice_integral}) gives an ice volume of \qty{168.2}{\micro\liter\per\gram} (\qty{0.5}{\kelvin\per\minute}) and \qty{154.8}{\micro\liter\per\gram} (\qty{2}{\kelvin\per\minute}), respectively.
However, the calculation of the ice mass is heavily influenced by the way this baseline of the heat flow $\dot{Q}_\text{B}$ is determined.
Since the temperature ranges to process the baseline correction are selected manually, results from this approach may differ depending on the person processing the data.
\zcref{fig:basle_line_ice_vol} illustrates the influence on the calculated ice volume, which can differ $\pm$\qty{10}{\micro\liter\per\gram} due to different selected temperature ranges to obtain the baseline function.
Therefore, the temperature ranges used for this correction were held constant for all further calculations.
This keeps different measurements, run with the same heating rate, comparable.

% The cumulative ice mass calculated using both methods is shown in \zcref{fig:ltdsc_result_heating_rate_b}.
% As can be seen in the zoomed in section, the manual approach gives more plausible curves, starting at roughly \qty{-47}{\celsius}, while the automatic approach seems to loose ice mass until \num{-38} ore \qty{-34}{\celsius}.
% This is caused by the intersection point of the base line corrected raw data with the x-axis.

While a slower heating rate causes improved results, the time required for a single measurement was deemed unpractical.
%It would not have been possible to obtain three measurements for every hydration time without a significant change in sample age between repeats.
Since the specimens are still wet and hydrating, the repeated measurements had to be carried out in a timely manner.
Therefore, it has been decided to use a faster heating rate of \qty{2}{\kelvin\per\minute}. % and the automatic approach, which should cause a comparable error within the same measurement regime.

\begin{figure}[htb]
    \centering
	%\tikzsetnextfilename{ltdsc_result_C2S_melt_bar} % Gibt den Dateinamen der externen Grafik vor
	\begin{subfigure}[t]{.49\textwidth}
	  \begin{tikzpicture}
		\begin{axis}[
			width  = \textwidth,
			height = 8cm,
			major x tick style = transparent,
			ybar=2*\pgflinewidth,
			ymin=0, ymax=370,
			bar width=12pt,
			ymajorgrids = true,
			ylabel = {cum. ice volume in \unit{\micro\liter\per\gram}},
			xticklabels={
				\qty{7}{\day},
				\qty{14}{\day},
				\qty{21}{\day},
				\qty{28}{\day}
			},
			xtick = data,
			scaled y ticks = false,
			enlarge x limits=0.15,
			ymin=0,
			legend cell align=left,
			legend style={
					at={(0.98,0.97)},
					anchor=north east,
					column sep=4pt
			},
			legend columns=-1,
			scaled y ticks = false,
			error bars/y dir=both,
			error bars/y explicit,
			ytick={0,100,200,300,400},
			visualization depends on={\thisrow{pores_b_std} \as \offset},
			nodes near coords,
			node near coords style={
				shift={(axis direction cs:0,\offset)},
				rotate=90,
				anchor=west,
				/pgf/number format/.cd,fixed zerofill,precision=1,
			},
		]
			\addplot[
				orange,fill=orange,mark=none,
				error bars/.cd,
				error bar style={black},
			] table [ x expr=\coordindex, y=pores_a_mean, y error=pores_a_std, col sep=comma] {figures/ltdsc/C2S_ice_vol_dev.csv};

			\addplot[
				cyan,fill=cyan,mark=none,
				error bars/.cd,
				error bar style={black}
			] table [ x expr=\coordindex, y=pores_b_mean, y error=pores_b_std, col sep=comma] {figures/ltdsc/C2S_ice_vol_dev.csv};

			% \addplot[
			% 	blue,fill=blue,mark=none,
			% 	error bars/.cd,
			% 	error bar style={black},
			% ] table [ x expr=\coordindex, y=pores_c_mean, y error=pores_c_std, col sep=comma] {figures/ltdsc/C2S_ice_vol_dev.csv};

			% \legend{gel pores, capillary pores, total pores}
		\end{axis}
	  \end{tikzpicture}

      \caption{Ice volume of \ac{C2S}}
 	  \label{fig:ltdsc_result_C2S_melt_bar}
	\end{subfigure}
	\hfill
	\begin{subfigure}[t]{.49\textwidth}
	  \begin{tikzpicture}
		\begin{axis}[
			width  = \textwidth,
			height = 8cm,
			major x tick style = transparent,
			ybar=2*\pgflinewidth,
			ymin=0, ymax=370,
			bar width=12pt,
			ymajorgrids = true,
			ylabel = {ice volume in \unit{\micro\liter\per\gram}},
			xticklabels={
				\qty{1}{\day},
				\qty{7}{\day},
				\qty{14}{\day},
				\qty{21}{\day},
				\qty{28}{\day}
			},
			xtick = data,
			enlarge x limits=0.15,
			ymin=0,
			legend cell align=left,
			legend style={
					at={(0.98,0.97)},
					anchor=north east,
					column sep=4pt
			},
			error bars/y dir=both,
			error bars/y explicit,
			legend columns=1,
			scaled y ticks = false,
			ytick={0,100,200,300,400},
			visualization depends on={\thisrow{pores_b_std} \as \offset},
			nodes near coords,
			node near coords style={
				shift={(axis direction cs:0,\offset)},
				rotate=90,
				anchor=west,
				/pgf/number format/.cd,fixed zerofill,precision=1,
			},
			yticklabel pos=right,
		]
			\addplot[
				orange,fill=orange,mark=none,
				error bars/.cd,
				error bar style={black},
			] table [ x expr=\coordindex, y=pores_a_mean, y error=pores_a_std, col sep=comma] {figures/ltdsc/C3S_ice_vol_dev.csv};

			\addplot[
				cyan,fill=cyan,mark=none,
				error bars/.cd,
				error bar style={black}
			] table [ x expr=\coordindex, y=pores_b_mean, y error=pores_b_std, col sep=comma] {figures/ltdsc/C3S_ice_vol_dev.csv};

			% \addplot[
			% 	blue,fill=blue,mark=none,
			% 	error bars/.cd,
			% 	error bar style={black},
			% ] table [ x expr=\coordindex, y=pores_c_mean, y error=pores_c_std, col sep=comma] {figures/ltdsc/C3S_ice_vol_dev.csv};

			\legend{gel pores, capillary pores}%, total pores}
		\end{axis}
	  \end{tikzpicture}
	  \caption{Ice volume of \ac{C3S}}
	  \label{fig:ltdsc_result_C3S_melt_bar}

	\end{subfigure}

    \caption{
		Ice volume per specimen weight for gel and capillary pores of \ac{C2S} and \ac{C3S} from \qtyrange{1}{28}{\day} of hydration.
		The data were calculated from three \ac{LTDSC}-measurements for each value.
		The baseline correction was done using temperature range from a lower range of \qtyrange{-21.5}{-20.5}{\celsius} (capillary pores) and \qtyrange{-57}{-56}{\celsius} (gel pores) and \qtyrange{9}{10}{\celsius} in the upper range.
	}
	\label{fig:ltdsc_result_melt_bar}
\end{figure}

As shown in \zcref{fig:ltdsc_result_C2S_melt_bar}, for \ac{C2S} the ice volume of the gel pores was calculated to be \qty{1.2}{\micro\liter\per\gram} at \qty{7}{\day} up to \qty{18.3}{\micro\liter\per\gram} after \qty{28}{\day} of hydration, while the capillary pores were reduced from \num{290} to \qty{233}{\micro\liter\per\gram} in the same time.
In contrast, \ac{C3S} has a higher ice volume accountable to the gel pores (\zcref{fig:ltdsc_result_C3S_melt_bar}).
The gel pore ice volume was increasing from \qty{6.6}{\micro\liter\per\gram} at an age of \qty{1}{\day} to \qty{30.3}{\micro\liter\per\gram} after \qty{28}{\day} of hydration.
In comparison to \ac{C2S}, less than half of the ice volume corresponding to capillary pores is detected in \ac{C3S} after \qty{7}{\day} of hydration.
The capillary pore space was reduced by approximately \qty{50}{\percent} from \qty{208}{\micro\liter\per\gram} (\qty{1}{\day}) to \qty{109}{\micro\liter\per\gram} (\qty{28}{\day}).
Furthermore, a significant slowdown of gel pore development and reduction in capillary pores can be seen in \ac{C3S} after \qty{14}{\day} of hydration.

As expected from the freezing curves, the overall ice volume and the ice volume correlated to the capillary pores is lower for \ac{C3S} compared to \ac{C2S}.
Nevertheless, both materials indicate a steady increase in gel pore and a reduction of capillary ice volume over time.
\zcref{fig:ltdsc_result_melt_correlation} indicates a roughly linear correlation of the increase in gel pore volume and the decrease in capillary pore volume until the \qty{28}{\day} mark.


\begin{figure}[htb]
    \centering

	\begin{tikzpicture}
		\begin{axis}[
			xlabel={gel pore volume in \unit{\micro\liter\per\gram}},
			ylabel={capillary pore volume in \unit{\micro\liter\per\gram}},
			xmin=0,xmax=35,
			ymin=0,ymax=350,
			width=.85\textwidth,
			height=6.5cm,
			grid=major,
			scatter/classes={
				1={mark=diamond*},
				7={mark=triangle*},
				14={mark=pentagon*},
				21={mark=*},
				28={mark=square*}
			},
			legend style={at={(1.02,0.02)}, anchor=south west},
            legend cell align={left},
		]

		% C3S linear fit
		\addplot [draw=none, forget plot] table [
			x=pores_a_mean, y=pores_b_mean, col sep=comma,
			y={create col/linear regression={y=pores_b_mean, x=pores_a_mean}}
		] {figures/ltdsc/C3S_ice_vol_dev.csv};
		\addplot [grey, thin, domain=0:30.3, forget plot] {\pgfplotstableregressiona*x + \pgfplotstableregressionb};

		% C2S linear fit
		\addplot [draw=none, forget plot] table [
			x=pores_a_mean, y=pores_b_mean, col sep=comma,
			y={create col/linear regression={y=pores_b_mean, x=pores_a_mean}}
		] {figures/ltdsc/C2S_ice_vol_dev.csv};
		\addplot [grey, thin, dashed, domain=0:18.3, forget plot] {\pgfplotstableregressiona*x + \pgfplotstableregressionb};

		% C3S
		\addplot [
			scatter, only marks, red, mark size=4pt,
			scatter src=explicit symbolic,
		] table [
			x=pores_a_mean, y=pores_b_mean, col sep=comma,
			meta=a
		] {figures/ltdsc/C3S_ice_vol_dev.csv};

		% C2S
		\addplot [
			scatter, only marks, orange, mark size=4pt,
			scatter src=explicit symbolic
		] table [
			x=pores_a_mean, y=pores_b_mean, col sep=comma,
			meta=a
		] {figures/ltdsc/C2S_ice_vol_dev.csv};

		\end{axis}

		% second axis for phase legend
		\begin{axis}[
			width=.85\textwidth,
			height=6.5cm,
			hide axis,
			xmin=0,xmax=1,ymin=0,ymax=1,
			legend style={at={(1.02,0.98)}, anchor=north west},
            legend cell align={left},
		]

		\addlegendimage{only marks, mark=square*, mark size=4pt, red}
		\addlegendentry{\ac{C2S}}
		\addlegendimage{only marks, mark=square*, mark size=4pt, orange}
		\addlegendentry{\ac{C3S}}

		\end{axis}

		% third axis for age legend
		\begin{axis}[
			width=.85\textwidth,
			height=6.5cm,
			hide axis,
			xmin=0,xmax=1,ymin=0,ymax=1,
			legend style={at={(1.02,0.02)}, anchor=south west},
            legend cell align={left},
		]

		\addlegendimage{only marks, mark=diamond*, mark size=4pt, grey}
		\addlegendentry{\qty{1}{\day}}
		\addlegendimage{only marks, mark=triangle*, mark size=4pt, grey}
		\addlegendentry{\qty{7}{\day}}
		\addlegendimage{only marks, mark=pentagon*, mark size=4pt, grey}
		\addlegendentry{\qty{14}{\day}}
		\addlegendimage{only marks, mark=*, mark size=4pt, grey}
		\addlegendentry{\qty{21}{\day}}
		\addlegendimage{only marks, mark=square*, mark size=4pt, grey}
		\addlegendentry{\qty{28}{\day}}

		\end{axis}
	\end{tikzpicture}

    \caption{
		Correlation between gel and capillary pore volume in \ac{C3S} and \ac{C2S} up to a hydration time of \qty{28}{\day}.
	}
	\label{fig:ltdsc_result_melt_correlation}
\end{figure}

These observations fit well to the microstructural development visible in \ac{SEM}-\ac{BSE} images of both materials, where the capillary porosity is reduced, since the hydrate phases grow into these larger pore spaces, causing new gel porosity (see \zcref{fig:hydration-rim-results}).

% I have decided to remove this, but leave it as a comment as a reminder, if the question arises.
% \begin{itemize}
% 	\color{red}
% 	\item difficult to calculate the pore volume fraction (in contrast to the pore volume per mass) % while this concern is true, since the actual volume of the specimen measured is unknown, I think it does not really matter, since this also was not done for the image analysis and is also unusual for MIP
% 	\item volume fraction of hydrate phase and unhydrated phase not fully known
% 	\item can be assumed using Area fraction from SEM images as described in \cite{kleiner.2024a} and \zcref{eq:}.
% \end{itemize}

\FloatBarrier
\subsubsection[MIP]{\acs*{MIP}}

\zcref{fig:mip_c2s_c3s} illustrates the \ac{MIP} results for \ac{C2S} and \ac{C3S}, respectively, after up to \qty{28}{\day} of hydration in comparison to the second freezing curves of the \ac{LTDSC}-measurements.
For these diagrams, the freeze curves were baseline corrected, which was performed by by calculating the linear regression of the data between \numrange{-55}{-50} and \qtyrange{-27}{-23}{\celsius} and subtracting this regression function from the mean data as shown in \zcref{fig:ltdsc_result_c2s_melt}.

\begin{figure}[!htb]
    \centering

    \begin{subfigure}[t]{\linewidth}
	  \begin{tikzpicture}
		\begin{semilogxaxis}[
			legend style={
				at={(.02,.96)},
				anchor=north west,
			},
			legend cell align={left},
			legend columns=2,
			xmin=1.5, xmax=5000,
			ymin=-.12, ymax=0,
			y dir=reverse,
			width=.83\linewidth,
			height=6cm,
			grid=major,
			xlabel={temperature in \unit{\celsius}},
			xlabel style={blue},
			xticklabel style={blue, font=\footnotesize},
			xticklabels={\num{-51.2}, \num{-25.3}, \num{-12.4}, \num{-4.6}, \num{-2.0}},
			xtick={2.5,5,10,25,50},
			ylabel={heat flow $\dot{Q}$ in \unit{\joule\per\second\per\gram}},
			scaled y ticks=false,
			scaled x ticks=false,
			axis x line*=top,
			axis y line*=left,
			ytick={0,-0.05,-0.1,-0.15,-0.2},
			yticklabel style={
				/pgf/number format/fixed,
				/pgf/number format/precision=2,
				/pgf/number format/fixed zerofill,
			},
		]

			\addplot[dotted, thick, samples=2, darkgrey,forget plot] coordinates {(3.19,-.2)(3.19,.2)};
			\node[darkgrey, rotate=90, anchor=west] at (axis cs: 2.50,-.035) {\scriptsize <\,\qty{3.19}{\nano\meter}};

			\addplot [ no marks, orange, forget plot, opacity=0.5] table [each nth point=4, filter discard warning=false, unbounded coords=discard,x index=3, y index=4, col sep=comma] {figures/ltdsc/C2S_7d_mean_freeze_segment.csv};
			\addplot [ no marks, red, forget plot, opacity=0.5] table [each nth point=4, filter discard warning=false, unbounded coords=discard,x index=3, y index=4, col sep=comma] {figures/ltdsc/C2S_14d_mean_freeze_segment.csv};
			\addplot [ no marks, green, forget plot, opacity=0.5] table [each nth point=4, filter discard warning=false, unbounded coords=discard,x index=3, y index=4, col sep=comma] {figures/ltdsc/C2S_21d_mean_freeze_segment.csv};
			\addplot [ no marks, blue, forget plot, opacity=0.5] table [each nth point=4, filter discard warning=false, unbounded coords=discard,x index=3, y index=4, col sep=comma] {figures/ltdsc/C2S_28d_mean_freeze_segment.csv};

			% the color legend is universal for all experiments. The linestyle in addlegendentry however, could not be changed for some reason.
            \addplot [ no marks, orange, thick ]  coordinates {(0.1,0) (.9,0)};
            \addlegendentry{\qty{7}{\day}};
            \addplot [ no marks, red, thick ]  coordinates {(0.1,0) (.9,0)};
            \addlegendentry{\qty{14}{\day}};
            \addplot [ no marks, green, thick ]  coordinates {(0.1,0) (.9,0)};
            \addlegendentry{\qty{21}{\day}};
            \addplot [ no marks, blue, thick ]  coordinates {(0.1,0) (.9,0)};
            \addlegendentry{\qty{28}{\day}};

		\end{semilogxaxis}
		\begin{semilogxaxis}[
			legend style={
				at={(.98,.96)},
				anchor=north east,
			},
			legend columns=-1,
			legend cell align={right},
			xmin=1.5, xmax=5000,
			ymin=0, ymax=1.6,
			width=.83\linewidth,
			height=6cm,
			grid=none,
			xlabel={pore radius $r_\text{P}$ in \unit{\nano\meter}},
			xlabel style={black},
			xticklabel style={black},
			ylabel={$dV/dlogR$\\pore volume in \unit{\milli\liter\per\gram}},
			ylabel style={align=center, darkgreen},
			yticklabel style={
				darkgreen,
				/pgf/number format/fixed,
				/pgf/number format/precision=1,
				/pgf/number format/fixed zerofill,
			},
			grid style={line width=.1pt, draw=gray!10},
			xticklabels={1,2.5,5,10,25,50, 100, 1000, 10000},
			xtick={1,2.5,5,10,25,50, 100, 1000, 10000},
			ytick={},
			scaled y ticks=false,
			scaled x ticks=false,
			axis y line*=right,
			axis x line*=bottom,
			log basis x=10, % seems to be necessary for "restrict x to domain" to work.
		]
			\addplot [ no marks, red, forget plot, dotted] table [skip first n=1, y index=4, x expr={1000*\thisrowno{1}}, col sep=semicolon] {figures/mip/HG C2S 14d.CSV};
			\addplot [ no marks, blue, forget plot, dotted] table [skip first n=1, y index=4, x expr={1000*\thisrowno{1}}, col sep=semicolon] {figures/mip/HG C2S 28d.CSV};


			%the domain restriction is somewhat weird. It should be 10:10000, but that does not work. 1:10000 at least gave the correct result.
			\addplot [ thick, no marks, red, forget plot, restrict x to domain=1:10000] table [skip first n=1, y index=4, x expr={1000*\thisrowno{1}}, col sep=semicolon] {figures/mip/HG C2S 14d.CSV};
			\addplot [ thick, no marks, blue, forget plot, restrict x to domain=1:10000] table [skip first n=1, y index=4, x expr={1000*\thisrowno{1}}, col sep=semicolon] {figures/mip/HG C2S 28d.CSV};

			% distinguish the experiments in the legend
			\addplot [ no marks, black, opacity=0.5] coordinates {(0.1,0) (.9,0)};
			\addlegendentry{\acs*{LTDSC}};
			\addplot [ thick, no marks, black, solid] coordinates {(0.1,0) (.9,0)};
			\addlegendentry{\acs*{MIP}};

		\end{semilogxaxis}
	  \end{tikzpicture}
	  \caption{\ac{C2S} after \num{14} and \qty{28}{\day} of hydration.}
	  \label{fig:mip_c2s}
    \end{subfigure}

    \begin{subfigure}[t]{\linewidth}
% 	\caption{
% 		\ac{MIP}-analysis of \ac{C2S} after \num{14} and \qty{28}{\day} of hydration ({\color{orange}right axis}) compared with selected \ac{LTDSC}-curves (left axis).
% 		Note that pore sizes below \qty{3.19}{\nano\meter} are undefined as noted in \zcref{eq:brun_freeze}.}
%     \label{fig:mip_c2s}
% \end{figure}

% \begin{figure}[!htb]
%     \centering
	  \begin{tikzpicture}
		\begin{semilogxaxis}[
				legend style={
					at={(.02,.96)},
					anchor=north west,
				},
				legend cell align={left},
				legend columns=2,
				xmin=1.5, xmax=5000,
				ymin=-.12, ymax=0,
				y dir=reverse,
				width=.83\linewidth,
				height=6cm,
				grid=major,
				xlabel={temperature in \unit{\celsius}},
				xlabel style={blue},
				xticklabel style={blue, font=\footnotesize},
				xticklabels={\num{-51.2}, \num{-25.3}, \num{-12.4}, \num{-4.6}, \num{-2.0}},
				xtick={2.5,5,10,25,50},
				ylabel={heat flow $\dot{Q}$ in \unit{\joule\per\second\per\gram}},
				scaled y ticks=false,
				scaled x ticks=false,
				axis x line*=top,
				axis y line*=left,
				ytick={0,-0.05,-0.1,-0.15,-0.2},
				yticklabel style={
					/pgf/number format/fixed,
					/pgf/number format/precision=2,
					/pgf/number format/fixed zerofill,
				},
			]

			\addplot[dotted, thick, samples=2, darkgrey,forget plot] coordinates {(3.19,-.2)(3.19,.2)};
			\node[darkgrey, rotate=90, anchor=west] at (axis cs: 2.50,-.020) {\scriptsize <\,\qty{3.19}{\nano\meter}};

			\addplot [ no marks, yellow, opacity=0.5, forget plot] table [each nth point=4, filter discard warning=false, unbounded coords=discard,x index=3, y index=4, col sep=comma] {figures/ltdsc/C3S_1d_mean_freeze_segment.csv};
            \addplot [ no marks, orange, opacity=0.5, forget plot] table [each nth point=4, filter discard warning=false, unbounded coords=discard,x index=3, y index=4, col sep=comma] {figures/ltdsc/C3S_7d_mean_freeze_segment.csv};
            \addplot [ no marks, red, opacity=0.5, forget plot] table [each nth point=4, filter discard warning=false, unbounded coords=discard,x index=3, y index=4, col sep=comma] {figures/ltdsc/C3S_14d_mean_freeze_segment.csv};
            \addplot [ no marks, green, opacity=0.5, forget plot] table [each nth point=4, filter discard warning=false, unbounded coords=discard,x index=3, y index=4, col sep=comma] {figures/ltdsc/C3S_21d_mean_freeze_segment.csv};
            \addplot [ no marks, blue, opacity=0.5, forget plot] table [each nth point=4, filter discard warning=false, unbounded coords=discard,x index=3, y index=4, col sep=comma] {figures/ltdsc/C3S_28d_mean_freeze_segment.csv};

			% the color legend is universal for all experiments. The linestyle in addlegendentry however, could not be changed for some reason.
            \addplot [ no marks, yellow, thick ]  coordinates {(0.1,0) (.9,0)};
            \addlegendentry{\qty{1}{\day}};
            \addplot [ no marks, orange, thick ]  coordinates {(0.1,0) (.9,0)};
            \addlegendentry{\qty{7}{\day}};
            \addplot [ no marks, red, thick ]  coordinates {(0.1,0) (.9,0)};
            \addlegendentry{\qty{14}{\day}};
            \addplot [ no marks, green, thick ]  coordinates {(0.1,0) (.9,0)};
            \addlegendentry{\qty{21}{\day}};
            \addplot [ no marks, blue, thick ]  coordinates {(0.1,0) (.9,0)};
            \addlegendentry{\qty{28}{\day}};

		\end{semilogxaxis}
		\begin{semilogxaxis}[
			legend style={
				at={(.98,.96)},
				anchor=north east,
			},
			legend columns=2,
			legend cell align={right},
			xmin=1.5, xmax=5000,
			ymin=0, ymax=.45,
			width=.83\linewidth,
			height=6cm,
			grid=none,
			xlabel={pore radius $r_\text{P}$ \unit{\nano\meter}},
			ylabel={$dV/dlogR$\\pore volume in \unit{\milli\liter\per\gram}},
			ylabel style={darkgreen, align=center},
			grid style={line width=.1pt, draw=gray!10},
			yticklabel style={
				darkgreen,
				/pgf/number format/fixed,
				/pgf/number format/precision=1,
				/pgf/number format/fixed zerofill,
			},
			xticklabels={1,2.5,5,10,25,50, 100, 1000, 10000},
			xtick={1,2.5,5,10,25,50, 100, 1000, 10000},
			ytick={},
			scaled y ticks=false,
			scaled x ticks=false,
			axis y line*=right,
			axis x line*=bottom,
			log basis x=10, % seems to be necessary for "restrict x to domain" to work.
		]
			\addplot [ orange, dotted, forget plot] table [skip first n=1, y index=4, x expr={1000*\thisrowno{1}}, col sep=semicolon] {figures/mip/HG C3S 7d.CSV};
			\addplot [ no marks, blue, dotted, forget plot] table [ skip first n=1, y index=4, x expr={1000*\thisrowno{1}}, col sep=semicolon] {figures/mip/HG C3S 28d.CSV};

			%the domain restriction is somewhat weird. It should be 10:10000, but that does not work. 1:10000 at least gave the correct result.
			\addplot [ thick, no marks, orange, forget plot, restrict x to domain=1:10000] table [skip first n=1, y index=4, x expr={1000*\thisrowno{1}}, col sep=semicolon] {figures/mip/HG C3S 7d.CSV};
			\addplot [ thick, no marks, blue,forget plot, restrict x to domain=1:10000] table [skip first n=1, y index=4, x expr={1000*\thisrowno{1}}, col sep=semicolon] {figures/mip/HG C3S 28d.CSV};

			% distinguish the experiments in the legend
			\addplot [ no marks, black, opacity=0.5] coordinates {(0.1,0) (.9,0)};
			\addlegendentry{\acs{LTDSC}};
			\addplot [ thick, no marks, black, solid] coordinates {(0.1,0) (.9,0)};
			\addlegendentry{\acs{MIP}};

		\end{semilogxaxis}
	  \end{tikzpicture}
	  \caption{\ac{C3S} after \num{7} and \qty{28}{\day} of hydration.}
	  \label{fig:mip_c3s}
    \end{subfigure}%
	\caption{
		\ac{MIP} analysis of \ac{C2S} after \num{14} and \qty{28}{\day} and \ac{C3S} after \num{7} and \qty{28}{\day} of hydration ({\color{darkgreen}right axis}) compared with selected \ac{LTDSC}-curves (left axis).
		Note that pore sizes below \qty{3.19}{\nano\meter} are undefined as noted in \zcref{eq:brun_freeze}.
	}
    \label{fig:mip_c2s_c3s}
\end{figure}

The \ac{LTDSC} curves above \qty{-12.4}{\celsius} ($r_\text{P}$=\qty{10}{\nano\meter}) are not reliable due to system inertia.
This is because the rapid freezing rate shifts the capillary pore space and larger pores become indistinguishable.
However, \ac{MIP} is considered fairly reliable above \qty{10}{\nano\meter}, while smaller pores are hardly detectable.
Therefore, both methods cover different pore space ranges.
\ac{LTDSC} indicates the development of the gel pores, while \ac{MIP} curves show the development of the capillary and larger pores.

The \ac{MIP} curves show a significant reduction in the pore size distribution for both materials from the first reliable measurable age (\ac{C2S}: \qty{14}{\day}; \ac{C3S}: \qty{7}{\day}) up to \qty{28}{\day} of hydration.
As already indicated by \ac{LTDSC}, the capillary pore space of \ac{C2S} is significantly higher than that of \ac{C3S} (note the scale difference of both diagrams).

\FloatBarrier
\subsection{Conclusions}

The characterisation of porosity in cementitious materials necessitates a multi-method approach to cover the full range of pore sizes and to validate the results.
The experiments using porous glass standards yielded promising results, serving as a successful benchmark for the \ac{LTDSC} method.
It was demonstrated that \ac{LTDSC} is capable of distinguishing between different pore sizes, validating its applicability for the subsequent analysis of cement-based materials.
Furthermore, the \ac{SEM}-based pore size analysis is highly sensitive to segmentation parameters, such as the blocksize of the adaptive thresholding algorithm and the watershed minimum-distance parameter.

For hydrated \ac{C3S} and \ac{C2S}, the combination of \ac{SEM}, \ac{FIBnt}, \ac{LTDSC}, and \ac{MIP} provided a comprehensive view of the pore structure development.
High-resolution \ac{SEM} and \ac{FIBnt} imaging of \ac{ArBIB}-prepared cross-sections allowed for the direct observation and quantification of the nano-porosity within \ac{CSHi}.
Both 2D and 3D image analysis yielded comparable pore size distributions with a dominant pore radius in the range of \qtyrange{9}{15}{\nano\meter}.
Comparison of \ac{ArBIB} preparation at room temperature and cryogenic conditions showed only minor differences in the measured pore structure, suggesting that thermal damage during preparation is small.
However, larger image areas of polished \ac*{CSHi} would be required to verify this assumption.

\ac{LTDSC} and \ac{MIP} measurements provided complementary volumetric data.
While \ac{MIP} is mostly limited to capillary pores (>\,\qty{10}{\nano\meter}), \ac{LTDSC} successfully distinguished between capillary and gel porosity.
However, the assignment to an absolute pore size using the approximation by \textcite{brun.1977} showed limitations, as the water in small gel pores freezes significantly below the boundary temperature of \qty{-40}{\degreeCelsius}.
Using the \textcite{brun.1977} relation indicated the presence of gel pores with derived pore sizes of $\approx$ \qty{2}{\nano\meter}, which differed from the image-based findings.
This is likely caused by the aforementioned limitations of the thermodynamic model at low temperatures.
However, the evaluation of the melting enthalpy in \ac{LTDSC} allowed for a quantitative assessment of pore volumes.

Comparing the methods regarding the hydration of \ac{C3S} and \ac{C2S}, these techniques consistently indicated a densification of the microstructure over time.
\ac{LTDSC} revealed a roughly linear correlation over time between the formation of gel porosity and the reduction of capillary porosity.
Specifically, \ac{C3S} exhibited a significantly higher volume of gel pores and a lower volume of capillary pores compared to \ac{C2S} at the same hydration age.
This aligns with the faster reaction kinetics of \ac{C3S} observed in other experiments.
\ac{MIP} confirmed the higher capillary porosity in \ac{C2S} pastes.
Consequently, while no single method can fully characterise the entire pore space, the combination of direct imaging (\ac{ArBIB} combined with \ac{SEM} and \ac{FIBnt}) give a good idea for the gel pore morphology within \ac*{CSHi}.

More indirect methods like \ac{LTDSC} and \ac{MIP} for volumetric quantification provide a robust description of the evolving pore structure, which is in accordance with the results based on image analysis shown in \zcref{chap:hydrate-rim}.
% TODO:  in \zcref{chap:hydrate-rim}, try to perform some kind of CLD image analysis of the porosity, remove large pores using a filter?